\textsc{The goal of this chapter is}
to generalise the following theorem:

\begin{theoremalph}[{\cite[Theorem~1]{ostrik03:modul-hopf}}]\label{finiterecon}
  Let \(\cat{C}\) be a finite tensor category
  and let \(\cat{M}\) be a finite abelian \(\cat{C}\)\hyp{}module category,
  such that the evaluation functor \(\blank \triangleright m \from \cat{C} \to \cat{M}\) is exact,
  for all \(m \in \cat{M}\).
  Then there exists an algebra object \(A \in \cat{C}\)
  such that there is an equivalence of \(\cat{C}\)\hyp{}module categories \(\modc{A} \simeq \cat{M}\).
\end{theoremalph}

\noindent{}This theorem makes several assumptions on \(\cat{C}\) and \(\cat{M}\):
\begin{itemize}[noitemsep]
  \item Finiteness assumptions: \(\cat{C}\) and \(\cat{M}\) are required to be finite abelian.
  \item The monoidal structure \(\blank \otimes \bblank\) of \(\cat{C}\)
  and the action functor \(\blank \triangleright \bblank\) of \(\cat{M}\)
  are both assumed to be exact in both variables.
  \item Rigidity assumptions on \(\cat{C}\);
  \ie, that all of its objects admit left and right duals with respect to its monoidal structure.
\end{itemize}

We will generalise \cref{finiterecon} in a way that greatly relaxes the first two kinds of assumptions,
and removes the third one altogether.
As mentioned before,
by~\cite[Example~2.20]{douglas19},
this reconstruction cannot solely rely on algebra objects in \(\cat{C}\),
and we have to approach \cref{finiterecon} from a more monadic point of view.
Our main result will be the following,
and presents a certain \emph{classification} of right exact lax \(\cat{C}\)\hyp{}module monads.

\begin{reptheorem}{mainnonsenseprojective,mainnonsenseinjective,projcorrespondence}
  Assume that \(\cat{C}\) and \(\cat{M}\) have enough projectives (injectives) and that there is an object \(\ell \in \cat{M}\) such that:
  \begin{itemize}[noitemsep]
    \item there is a right adjoint \(\hom{\ell,\blank}\) (left adjoint \(\cohom{\ell,\blank}\)) to \(\blank \triangleright \ell\);
    \item for \(x \in \cat{C}\) projective (injective), the object \(x \triangleright \ell\) is projective (injective);
    \item any projective (injective) object of \(\cat{M}\) is a direct summand of an object of the form \(x \triangleright \ell\), for \(x\) projective (injective).
  \end{itemize}
  Let \(T\) be the monad \(\hom{\ell, \blank\,\triangleright\,\ell}\).
  Then \(\cat{M} \simeq \cat{C}^T\), and the \(\cat{C}\)\hyp{}module structure of the category of \(T\)\hyp{}modules is extended from the Kleisli category. Furthermore, this extends to a bijection
  \begin{align*}
    \faktor{\{\,(\cat{M},\ell) \text{ as above}\,\}}{(\cat{M} \simeq \cat{N})}
    &\ \leftrightarrow\,
      \Bigg\{
      \begin{aligned}
        &\text{Right exact lax \(\cat{C}\)\hyp{}module}\\
        &\text{monads on \(\cat{C}\)}
      \end{aligned}
      \Bigg\}
      /\, (\cat{C}^T \simeq \cat{C}^S) \\
    (\cat{M},\ell) &\longmapsto \hom{\ell,-\lact \ell} \\
    (\cat{C}^T, T1) &\longmapsfrom T
  \end{align*}
\end{reptheorem}

In the case of tensor categories,
this statement may on the surface seem inapplicable:
a tensor category that is not finite will generally not have enough projectives or injectives.
However, since the category underlying a tensor category is equivalent to the category of finite-dimensional comodules for a coalgebra,
the ind-completion of a tensor category has enough injectives.
Further, since multitensor categories can be realised as categories of compact objects in locally finitely presentable categories,
we may use an appropriate variant of the special adjoint functor theorem, see~\cref{saftrex,saftlex},
characterising when the internal cohom exists for an object in \(\Ind(\cat{M})\).
As such, we are able to ``pull back'' the appropriate versions of
\cref{thm:main-result-rigid-projective-case,thm:main-result-rigid-injective-case}
to the underlying categories.

\begin{reptheorem}{thm:multitensor-ind-reconstruction}
  Let \(\cat{C}\) be a multitensor category,
  and let \(\cat{M}\) be an abelian \(\cat{C}\)\hyp{}module category
  such that the \(\Ind(\cat{C})\)\hyp{}module category \(\Ind(\cat{M})\)
  admits a coclosed \(\Ind(\cat{C})\)\hyp{}injective \(\Ind(\cat{C})\)\hyp{}cogenerator.

  Then there is a coalgebra object \(C\) in \(\Ind(\cat{C})\)
  such that \(\Ind(\cat{M}) \simeq \mathsf{Comod}_{\Ind(\cat{C})}C\).
  Further, \(\cat{M}\) is the category of compact \(C\)\hyp{}comodule objects..
\end{reptheorem}

Consider the special case where a tensor category \(\cat{C}\) admits a fibre functor
to the category \(\vect\) of finite-dimensional vector spaces,
and is thus monoidally equivalent to the category \(\tetramodfd[H]\) of finite-dimensional comodules
over a finite-dimensional Hopf algebra \(H\), see~\cite{ulbrich1990:HopfAlgebrasRigid}.
\Cref{finiterecon} realises a finite \(\cat{C}\)\hyp{}module category \(\cat{M}\) as
the category of finite-dimensional modules over a finite-dimensional \(H\)\hyp{}comodule algebra.
Since \(H\) and \(A\) are finite-dimensional,
the category of \(A\)\hyp{}modules is equivalent
as an \(\tetramodfd[H]\)\hyp{}module category
to the category of comodules for the \(H\)\hyp{}comodule coalgebra \(A^{\ast}\).
In this sense, the following Hopf-theoretic corollary of the above theorem about ind-completions is an immediate generalisation of \cref{finiterecon}.

\begin{reptheorem}{infinitecomodulecoalgebras}
  Let \(H\) be a Hopf algebra and let \(\cat{C} = \tetramodfd[H]{}\),
  meaning that there is a monoidal equivalence \(\Ind(\cat{C}) \simeq \Tetramod[H]{}\).
  Let \(\cat{M}\) be an abelian \(\cat{C}\)\hyp{}module category
  such that \(\Ind(\cat{M})\)
  admits a coclosed \(\Ind(\cat{C})\)\hyp{}injective \(\Ind(\cat{C})\)\hyp{}cogenerator.

  Then there exists a (possibly infinite-dimensional) \(H\)\hyp{}comodule coalgebra \(C\),
  such that there is an equivalence \(\Ind(\cat{M}) \simeq \mathrm{Comod}_{H}C\)
  of \(\Ind(\cat{C})\)\hyp{}module categories,
  restricting to a \(\cat{C}\)\hyp{}module equivalence \(\cat{M} \simeq \mathsf{comod}_{H}C\).
\end{reptheorem}

Since we do not assume rigidity for our most general theorems,
we instead invoke the bicategorical Yoneda lemma to realise the objects of \(\cat{C}\) as the \(\cat{C}\)\hyp{}module endofunctors of \(\cat{C}\) itself.
Then, we view the \emph{lax} \(\cat{C}\)\hyp{}module endofunctors of \(\cat{C}\) as one appropriate generalisation of objects in \(\cat{C}\),
and the lax \(\cat{C}\)\hyp{}module monads as the corresponding generalisation of algebra objects in \(\cat{C}\).
Using the reconstruction results of \cref{sec:comodule-monads},
the right adjoint of a (strong) \(\cat{C}\)\hyp{}module functor is canonically lax,
and thus the (co)monads we study are canonically (op)lax \(\cat{C}\)\hyp{}module functors.

\section{Extending module structures}

\textsc{An additional difficulty we encounter}
is that while the Kleisli category of a lax \(\cat{C}\)\hyp{}module monad is naturally a \(\cat{C}\)\hyp{}module category,
the same is not true for the Eilenberg--Moore category.
In this section, we shall investigate under which conditions the \(\cat{C}\)\hyp{}module structure of the former lifts essentially uniquely to the latter.
More precisely,
we first establish uniqueness in \cref{thm:one-module-structure-on-EM},
and then complement that with an existence result of such an extension in \cref{extendingalways},
under the assumption of right exactness (left exactness for comonads) of the (co)monad,
using so-called Linton coequalisers and multicategorical techniques similar to those in~\cite{aguiar18:monad}.

\begin{hypothesis}\label{hyp:deligne-reconstr}
  From now on until the end of this thesis, we implicitly assume all categories and functors to be \(\Bbbk\)\hyp{}linear,
  for a field \(\Bbbk\).
\end{hypothesis}

The fact that we only consider right exact lax \(\cat{C}\)\hyp{}module monads
is exemplified by the following fundamental fact.

\begin{proposition}[{\cite[Proposition~3.2]{ben-zvi18:integ}}]\label{prop:abelian-EM-cat}
  Let \(\cat{A}\) be abelian,
  \(T\) a right exact monad on \(\cat{A}\),
  and \(S\) a left exact comonad on \(\cat{A}\).
  Then \(\cat{A}^T\) and \(\cat{A}^S\) are abelian.
\end{proposition}

\begin{proposition}\label{locomonads}
  Let \(\cat{A}\) be a locally finite abelian category,
  and let \(S\) be a left exact,
  finitary comonad on \(\Ind(\cat{A})\).
  Then \({\Ind(\cat{A})}^{S}\) also is of the form \(\Ind(\cat{E})\),
  for a locally finite abelian category \(\cat{E}\).
  Further, \(\cat{E}\) can be chosen to be the category of compact objects in \({\Ind(\cat{A})}^S\),
  and it can be characterised as the objects sent to compact objects under the forgetful functor \(F^S\from {\Ind(\cat{A})}^S \to \Ind(\cat{A})\).

  In particular, if \(S\) is a left exact comonad on \(\cat{A}\),
  we have \({\Ind(\cat{A})}^{\Ind(S)} \simeq \Ind(\cat{A}^S)\)
  and \(\cat{A}^S\) is locally finite abelian.
\end{proposition}
\begin{proof}
  Let \(D\) be the \(\Bbbk\)\hyp{}coalgebra such that \(\cat{A} \simeq \tetramodfd[D]{}\).
  Then we also have that \(\Ind(\cat{A}) \simeq \Tetramod[D]{}\).
  In particular, \(S\) is a comonad on \(\Tetramod[D]{}\).
  By \cref{TakeuchiEilenbergWatts}, there is a \(D\)\hyp{}\(D\)\hyp{}bicomodule \(C\) such that \(S \cong C \cotens_D \blank\).

  Similarly to~\cite[Remark~2.4]{takeuchi77:morit},
  under this isomorphism the comonad structure on \(S\) corresponds
  to maps \(C \to C \square_{D} C\) and \(C \to D\),
  which endow \(C\) with the structure of a \(\Bbbk\)\hyp{}coalgebra,
  together with a coalgebra morphism \(C \to D\).
  Formally, a \(S\)\hyp{}comodule in \(\Tetramod[D]{}\) is a \(D\)\hyp{}comodule
  together with a \(C\)\hyp{}comodule structure
  that restricts to the given \(D\)\hyp{}comodule along the coalgebra morphism \(C \to D\).
  Morphisms of \(S\)\hyp{}comodules in \(\Tetramod[D]{}\) are precisely \(C\)\hyp{}comodule morphisms,
  and so we have
  \begin{equation}\label{eq:locomonads:equiv}
    {\Ind(\cat{A})}^{S} \simeq {(\Tetramod[D]{})}^{S} \simeq \Tetramod[C]{}.
  \end{equation}
  Thus we may set \(\cat{E} = \tetramodfd[C]{}\).
  The characterisation of compact objects in \({(\Tetramod[D]{})}^{S}\)
  in terms of the images of the functor \(F^{S}\) follows immediately from observing that
  under the equivalence of \cref{eq:locomonads:equiv},
  the functor \(F^{S}\) is simply the restriction functor \(\Tetramod[C]{} \to \Tetramod[D]{}\).

  The second part of the statement follows by observing that \(\cat{A}^S\) is the category of compact objects in \({\Ind(\cat{A})}^{\Ind(S)}\),
  by the first part of the statement,
  and by observing that \(\cat{A}^S\) is also the category of compact objects in \(\Ind(\cat{A}^S)\).
  Since both \(\Ind(\cat{A}^S)\) and \({\Ind(\cat{A})}^{\Ind(S)}\) are locally finitely presentable,
  the equivalence between their respective categories of compact objects establishes
  an equivalence between the categories themselves.
\end{proof}

\begin{remark}\label{rmk:em-iso-and-so-on}
  Suppose that \(T\) is a monad on the category \(\cat{C}\).
  For \(T\)\hyp{}modules \((x, \alpha)\) and \((y, \beta)\),
  notice that there is a bijection
  \begin{equation}\label{eq:fincorestr:iso}
    \begin{aligned}
      \cat{C}^T(Tx, y) &\,\ \cong\,\ \cat{C}(x, y) \\
      (f\from Tx \to y) &\mapsto (f \circ \eta_x) \\
      (\beta \circ Tg) &\longmapsfrom (g\from x \to y).
    \end{aligned}
  \end{equation}
  This induces an isomorphism \(\cat{C}^T(T(\blank), \bblank) \cong \cat{C}(\blank, \bblank)\).
  Using the definition of \(\iota\) from \cref{denotedbyiota},
  the left-hand side is also obtained from the functor
  \[
    \cat{C}^T
    \xrightarrow{\ \yo\ } [{(\cat{C}^T)}^{\op}, \kVect]
    \xrightarrow{\ \iota^{\ast}\ } [{\cat{C}_T}^{\op}, \kVect].
  \]
\end{remark}

The analogue of \cref{locomonads} for finite abelian categories is simpler.

\begin{proposition}\label{finitemonads}
  Let \(T\) be a right exact monad on a finite abelian category \(\cat{A}\).
  Then \(\cat{A}^T\) is a finite abelian category.

  Further, any projective object in \(\cat{A}^T\) is a direct summand of one of the form \(\iota(P)\),
  where \(\iota\from \cat{A}_T \to \cat{A}^T\) is the canonical embedding of \cref{denotedbyiota} and \(P \in \proj{\cat{A}}\).
  Denoting the category of objects of this form by \(\mathsf{Kl}_p(T)\),
  we obtain an equivalence
  \[
    \proj{\cat{A}^T} \simeq \overline{\mathsf{Kl}_{p}(T)}
  \]
  between the projectives of \(\cat{A}^T\) and the Cauchy completion of \(\mathsf{Kl}_p(T)\).
\end{proposition}
\begin{proof}
  Let \(A\) be a finite\hyp{}dimensional \(\Bbbk\)\hyp{}algebra such that \(\cat{A} \simeq \lmod{A}\).
  Then there is an \(A\)\hyp{}\(A\)\hyp{}bimodule \(B\) such that \(T \cong B \otimes_{A} \blank\).
  Similarly to the proof of \cref{locomonads},
  monad structures on \(T\) correspond to pairs consisting of a finite\hyp{}dimensional \(\Bbbk\)\hyp{}algebra structure on \(B\) and an algebra homomorphism \(A \to B\).
  Under this correspondence we have \(\cat{A}^T \simeq \lmod{B}\).
  The latter category is clearly a finite abelian category.

  Now suppose that \(P \in \proj{\cat{A}}\).
  Applying \cref{rmk:em-iso-and-so-on}, one obtains an isomorphism \(\cat{A}^T(\iota(P), \blank) \cong \cat{A}(P, \blank)\),
  proving the exactness of the left-hand side, and thus projectivity of \(\iota(P)\).
  For \((X, \alpha) \in \cat{A}^T\), using the fact that \(\cat{A}\) has enough projectives,
  we may fix an epimorphism \(q\from Q \sur X\) in \(\cat{A}\),
  where \(Q \in \proj{\cat{A}}\).
  For the latter claim, notice that there is a composite epimorphism
  \[
    \iota(P) \xtwoheadrightarrow{\iota(q)} \iota(X) = TX \xtwoheadrightarrow{\ \alpha\ } X,
  \]
  the first part of which is epic by right exactness of \(\iota\).
\end{proof}

\begin{lemma}\label{finitecorestriction}
  Let \(T\) be a monad on \(\cat{A}\).
  The embedding
  \(\cat{A}^T \longhookrightarrow [\cat{A}_T^{\op}, \kVect]\)
  can be corestricted to an embedding
  \(\cat{A}^T \longhookrightarrow \mathrm{Fin}_{\mathrm{co}}(\cat{A}_T)\) into the finite cocompletion of \(\cat{A}_T\).
\end{lemma}
\begin{proof}
  Let \((x, \alpha)\) and \((y, \beta)\) be two \(T\)\hyp{}algebras.
  Recall that the coequaliser
  \[
    \begin{mytikzcd}[ampersand replacement=\&]
      {T^{2}x} \& {Tx} \& x
      \arrow["{\mu_{x}}", shift left=1, from=1-1, to=1-2]
      \arrow["{T \alpha}"', shift right=1, from=1-1, to=1-2]
      \arrow["{\alpha}", twoheadrightarrow, from=1-2, to=1-3]
    \end{mytikzcd}
  \]
  in \(\cat{A}^T\) is an absolute coequaliser\note{%
    A coequaliser is called \emph{absolute} if it is preserved by any functor;
    see~\cite{pare69:absol} for details.%
  }
  in \(\cat{A}\);
  thus, its image under the Yoneda embedding \(\yo\) yields a coequaliser
  \[
    \begin{mytikzcd}[ampersand replacement=\&]
      {\cat{A}(\blank, T^2x)} \& {\cat{A}(\blank, Tx)} \& {\cat{A}(\blank, x).}
      \arrow["{{(\mu_x)}_{*}}", shift left=1, from=1-1, to=1-2]
      \arrow["{{(T \alpha)}_{*}}"', shift right=1, from=1-1, to=1-2]
      \arrow["{\alpha_{*}}", twoheadrightarrow, from=1-2, to=1-3]
    \end{mytikzcd}
  \]
  Passing under the isomorphism of \cref{eq:fincorestr:iso},
  one obtains a coequaliser
  \[
    \begin{mytikzcd}[ampersand replacement=\&]
      {\cat{A}^T(T(\blank), T^2x)} \& {\cat{A}^T(T(\blank), Tx)} \& {\cat{A}^T(T(\blank), x),}
      \arrow["{{(\mu_x)}_{*}}", shift left=1, from=1-1, to=1-2]
      \arrow["{{(T \alpha)}_{*}}"', shift right=1, from=1-1, to=1-2]
      \arrow["{\alpha_{*}}", twoheadrightarrow, from=1-2, to=1-3]
    \end{mytikzcd}
  \]
  which, in turn, is isomorphic to
  \[
    \begin{mytikzcd}[ampersand replacement=\&]
      {\cat{A}_T(\blank, Tx)} \& {\cat{A}_T(\blank, x)} \& {\cat{A}^T(T(\blank), x).}
      \arrow["{{(\mu_x)}_{*}}", shift left=1, from=1-1, to=1-2]
      \arrow["{{(T\alpha)}_{*}}"', shift right=1, from=1-1, to=1-2]
      \arrow["{\alpha_{*}}", twoheadrightarrow, from=1-2, to=1-3]
    \end{mytikzcd}
  \]
  This proves the result.
\end{proof}

\begin{proposition}\label{prop:badmagic}
  Let \(T\) be a right exact monad on an abelian category \(\cat{A}\).
  The inclusion functor \(\cat{A}^T \longhookrightarrow \mathrm{Fin}_{\mathrm{co}}(\cat{A}_T)\) has a left adjoint, and the counit of the adjunction is a natural isomorphism.
  Further, the left adjoint \(\mathrm{Fin}_{\mathrm{co}}(\cat{A}_T) \to \cat{A}^T\)
  is the right exact extension of the inclusion \(\iota \from \cat{A}_T \longhookrightarrow \cat{A}^T\).
\end{proposition}
\begin{proof}
  By \cref{prop:abelian-EM-cat},
  the \(\cat{A}^T\) is finitely cocomplete.
  Similarly to~\cite[Proposition~6.3.1]{kashiwara06:categ},
  there is a functor \(\mathrm{Fin}_{\mathrm{co}}(\cat{A}^T) \to \cat{A}^T\)
  that extends the functor \(\mathrm{Id}\from \cat{A}^T \to \cat{A}^T\),
  defined by sending some \(x = \widetilde{\colim_i}\,x_i\) in \(\mathrm{Fin}_{\mathrm{co}}(\cat{A}^T)\)
  to \(\colim_i x_i\) in \(\cat{A}^T\),
  where \(\widetilde{\colim_i}\) denotes the formally added colimit.
  The functor \(\mathrm{Fin}_{\mathrm{co}}(\cat{A}^T) \to \cat{A}^T\)
  is left adjoint to the inclusion \(\cat{A}^T \longhookrightarrow \mathrm{Fin}_{\mathrm{co}}(\cat{A}^T)\).

  On the other hand, there is an adjunction
  \begin{equation}\label{eq:kan-adjunction}
    \begin{mytikzcd}[ampersand replacement=\&]
      {[\cat{A}_T^{\op}, \kVect]} \& {[{(\cat{A}^T)}^{\op}, \kVect],}
      \arrow[""{name=0, anchor=center, inner sep=0}, "{\iota_{!}}", shift left=2, from=1-1, to=1-2]
      \arrow[""{name=1, anchor=center, inner sep=0}, "{\iota^{\ast}}", shift left=2, from=1-2, to=1-1]
      \arrow["\dashv"{anchor=center, rotate=-90}, draw=none, from=0, to=1]
    \end{mytikzcd}
  \end{equation}
  where \(\iota_{{!}}\) is the left Kan extension of \(\yo_{\cat{A}^T} \circ \iota\) along \(\yo_{\cat{A}_T}\),
  see for example~\cite[Corollary~3.3]{stroinski24:ident}.
  \Cref{eq:kan-adjunction} restricts to an adjunction
  \[
    \begin{mytikzcd}[ampersand replacement=\&]
      {\mathrm{Fin}_{\mathrm{co}}(\cat{A}_T)} \& {\mathrm{Fin}_{\mathrm{co}}(\cat{A}^T),}
      \arrow[""{name=0, anchor=center, inner sep=0}, "{\iota_{!}}", shift left=2, from=1-1, to=1-2]
      \arrow[""{name=1, anchor=center, inner sep=0}, "{\iota^{\ast}}", shift left=2, from=1-2, to=1-1]
      \arrow["\dashv"{anchor=center, rotate=-90}, draw=none, from=0, to=1]
    \end{mytikzcd}
  \]
  since, by definition, \(\iota_{!}\) sends finite colimits of representables to finite colimits of representables,
  and \(\iota^{\ast}\) has the same property, since it preserves colimits, and, by \cref{finitecorestriction}, sends representables to finite colimits of representables.
  Thus we have the following commutative diagram
  \[
    \begin{mytikzcd}[ampersand replacement=\&]
      {\cat{A}^T} \& {\mathrm{Fin}_{\mathrm{co}}(\cat{A}^T)} \& {\cat{A}^T} \\
      \& {\mathrm{Fin}_{\mathrm{co}}(\cat{A}_T)} \& {\cat{A}_T}
      \arrow[""{name=0, anchor=center, inner sep=0}, shift right=2, from=1-1, to=1-2]
      \arrow[curve={height=14pt}, hook', from=1-1, to=2-2]
      \arrow[""{name=1, anchor=center, inner sep=0}, shift right=2, from=1-2, to=1-1]
      \arrow[""{name=2, anchor=center, inner sep=0}, "{\iota^{\ast}}", shift left=2, from=1-2, to=2-2]
      \arrow[hook', from=1-3, to=1-2]
      \arrow[""{name=3, anchor=center, inner sep=0}, "{\iota_{!}}", shift left=2, from=2-2, to=1-2]
      \arrow["\iota"', from=2-3, to=1-3]
      \arrow[hook', from=2-3, to=2-2]
      \arrow["\dashv"{anchor=center, rotate=-90}, draw=none, from=1, to=0]
      \arrow["\dashv"{anchor=center}, draw=none, from=3, to=2]
    \end{mytikzcd}
  \]
  This realises the embedding \(\cat{A}^T \longhookrightarrow \mathrm{Fin}_{\mathrm{co}}(\cat{A}_T)\)
  as the composition of two right adjoints,
  and thus as a right adjoint.
  Lastly, the composite
  \[
    \cat{A}_T \to \cat{A}^T \to \mathrm{Fin}_{\mathrm{co}}(\cat{A}^T) \to \cat{A}^T
  \]
  is naturally isomorphic to \(\iota \from \cat{A}_T \longhookrightarrow \cat{A}^T\),
  which proves the latter statement.
\end{proof}

\begin{remark}
  In other words, \cref{prop:badmagic} says that
  for a right exact monad \(T\) the category \(\cat{A}^T\) is
  a \emph{reflective subcategory}—one where the inclusion functor has a left adjoint—%
  of \(\mathrm{Fin}_{\mathrm{co}}(\cat{A}_T)\).
\end{remark}

The next result proves the essential uniqueness of the module structure
on the Eilenberg--Moore category,
under the assumption that it is ``induced'' from the Kleisli category.

\begin{theorem}\label{thm:one-module-structure-on-EM}
  Let \(\cat{C}_{T}\) denote the Kleisli category for a monad \(T\) on \(\cat{C}\),
  equipped with a fixed left \(\cat{C}\)\hyp{}module structure.
  Equip \(\cat{C}^T\) with two left \(\cat{C}\)\hyp{}module category structures
  \(\cat{C}^T_{1}\) and \(\cat{C}^T_{2}\),
  such that the inclusion
  \(\iota \from \cat{C}_T \longhookrightarrow \cat{C}^T\)
  gives strong \(\cat{C}\)\hyp{}module functors
  \(\cat{C}_{T} \to \cat{C}^T_{1}\)
  and \(\cat{C}_{T} \to \cat{C}^T_{2}\).

  Then there is an equivalence \(\cat{C}^T_{1} \simeq \cat{C}^T_{2}\) of left \(\cat{C}\)\hyp{}module categories.
\end{theorem}
\begin{proof}
  Consider the following diagram:
  \[
    \begin{mytikzcd}[ampersand replacement=\&]
      \& {\cat{C}_{T}} \\
      \& {\mathrm{Fin}_{\mathrm{co}}(\cat{C}_{T})} \\
      {\cat{C}^T_{1}} \&\& {\cat{C}^T_{2}}
      \arrow[from=1-2, to=2-2]
      \arrow["\iota"', curve={height=18pt}, from=1-2, to=3-1]
      \arrow["\iota", curve={height=-18pt}, from=1-2, to=3-3]
      \arrow[""{name=0, anchor=center, inner sep=0}, "{\iota^1_{!}}", shorten <=6pt, shift left=2, from=2-2, to=3-1]
      \arrow[""{name=1, anchor=center, inner sep=0}, "{\iota^2_{!}}"', shorten <=6pt, shift right=2, from=2-2, to=3-3]
      \arrow[""{name=2, anchor=center, inner sep=0}, "{G_{1}}", shift left=2, shorten >=6pt, from=3-1, to=2-2]
      \arrow[""{name=3, anchor=center, inner sep=0}, "{G_{2}}"', shift right=2, shorten >=6pt, from=3-3, to=2-2]
      \arrow["\dashv"{anchor=center, rotate=62}, draw=none, from=1, to=3]
      \arrow["\dashv"{anchor=center, rotate=118}, draw=none, from=0, to=2]
    \end{mytikzcd}
  \]
  In particular,
  by \cref{prop:dual-module-adjunctions},
  in this diagram every functor is a lax \(\cat{C}\)\hyp{}module functor,
  and every adjunction is a \(\cat{C}\)\hyp{}module adjunction.
  We claim that \(\iota^2_{!} G_{1}\from \cat{C}^T_{1} \to \cat{C}^T_2\) and \(\iota^1_{!} G_{2} \from \cat{C}^T_2 \to \cat{C}^T_1\)
  define mutually quasi-inverse equivalences of \(\cat{C}\)\hyp{}module categories.
  Indeed,
  \[
    \iota^1_{!} G_{2} \iota^2_{!} G_{1} \xrightarrow{\varepsilon_{2} \varepsilon_{1}} \mathrm{Id}_{\cat{C}^T_1}
  \]
  is a \(\cat{C}\)\hyp{}module isomorphism.

  Similarly, there is a \(\cat{C}\)\hyp{}module isomorphism \(\mathrm{Id}_{\cat{C}^T_2} \cong \iota^2_{!} G_{1} \iota^1_{!} G_{2}\).
\end{proof}

\subsection{Extendable monads}

\begin{definition}\label{extendable}
  We call a lax \(\cat{C}\)\hyp{}module monad \(T\) on a \(\cat{C}\)\hyp{}module category \(\cat{M}\) \emph{extendable}
  if the category \(\cat{M}^T\) is a \(\cat{C}\)\hyp{}module category,
  for which the embedding \(\iota\from \cat{M}_T \longhookrightarrow \cat{M}^T\) is a strong \(\cat{C}\)\hyp{}module functor.

  If \(T\) is extendable,
  we refer to the essentially unique \(\cat{C}\)\hyp{}module structure on \(\cat{M}^T\) coming from \(\cat{M}_T\)
  as the \emph{extended} \(\cat{C}\)\hyp{}module structure on \(\cat{M}^T\).
\end{definition}

Analogously to \cref{extendable},
one defines extendable oplax \(\cat{C}\)\hyp{}module comonads
and extended module structures on their comodules.

\begin{proposition}\label{prop:strong-module-monads-are-extendable}
  Any strong \(\cat{C}\)\hyp{}module monad \((T, \mu, \eta)\from \cat{M} \to \cat{M}\) is extendable.
\end{proposition}
\begin{proof}
  Since \(T\) is in particular a lax \(\cat{C}\)\hyp{}module monad,
  by \cref{strongcomparisons}, the comparison functor \(K_T\from \cat{M}_T \to \cat{M}\) is a strong \(\cat{C}\)\hyp{}module functor.
  Similarly, since \(T\) is in particular an oplax \(\cat{C}\)\hyp{}module monad,
  again by \cref{strongcomparisons}, we conclude that also the comparison functor \(K^T\from \cat{M} \to \cat{M}^T\) is a strong \(\cat{C}\)\hyp{}module functor.
  Thus, \(\iota = K^T K_T\) is a strong \(\cat{C}\)\hyp{}module functor.
\end{proof}

\begin{corollary}
  If \(\cat{C}\) is left rigid, then any lax \(\cat{C}\)\hyp{}module monad is extendable.
\end{corollary}
\begin{proof}
  By \cref{prop:lax-is-strong-when-rigid},
  a lax \(\cat{C}\)\hyp{}module monad is automatically a strong \(\cat{C}\)\hyp{}module monad.
  The result follows by \cref{prop:strong-module-monads-are-extendable}.
\end{proof}

\subsection{Semisimple monads}\label{sec:semisimple-monads}

\textsc{In this section we want to take a look}
at \emph{semisimple} monads—those whose category of modules is semisimple.
As it turns out, these monads are always extendable.
Recall the definition of the Cauchy completion of a category from \cref{sec: trep-rigid-fgp}.
The following result follows immediately from the fact that
for a monad \(T\from \cat{C}\to \cat{C}\),
the forgetful functor \(U^T\from \cat{C}^T \to \cat{C}\) creates limits.

\begin{lemma}
  The Eilenberg--Moore category of a monad on a Cauchy complete category is itself Cauchy complete.
\end{lemma}

\begin{proposition}\label{prop:SemisimpleEM}
  Let \(T\from \cat{A} \to \cat{A}\) be a right exact monad on an abelian category \(\cat{A}\).
  The category \(\cat{A}^T\) is semisimple if and only if
  \(\overline{\cat{A}_T}\) is.
  In this semisimple setting, the extension \(\overline{\iota}\from \overline{\cat{A}_T} \to \cat{A}^T\)
  of the canonical embedding \(\iota\from \cat{A}_T \to \cat{A}^T\) from \cref{denotedbyiota}
  to the Cauchy completion is an equivalence of categories.
\end{proposition}
\begin{proof}
  Since \(\cat{A}\) is Cauchy complete, the adjunction \(\adj{F_{T}}{U_T}{\cat{A}}{\cat{A}_T}\)
  extends to an adjunction \(\adj{\overline{F_{T}}}{\overline{U_{T}}}{\cat{A}}{\overline{\cat{A}_{T}}}\),
  such that \(\overline{U_T} \overline{F_T} = T\).

  It is easy to see that the resulting comparison functor \(K\from \overline{\cat{A}_T} \to \cat{A}^T\) is naturally isomorphic to the extension \(\overline{\iota}\) of \(\iota\).
  Further, since \(\iota\) is fully faithful, so is \(\overline{\iota}\), hence it reflects zero objects and so do \(U_T = U^T \circ \iota\)
  and \(\overline{U_T}\).

  Assume that \(\cat{A}^T\) is semisimple.
  The monomorphisms and epimorphisms in \(\cat{A}^T\) are split by \cref{toomuchTV},
  and hence they are reflected under the fully faithful functor \(\overline{\iota}\).
  Thus, \(\overline{\cat{A}_T}\) is an abelian category and \(\overline{\iota}\) is an exact functor.
  In particular \(\overline{U_T} \cong U_T \circ \overline{\iota}\) is exact and faithful,
  so it reflects zero objects.
  By \cref{thm:SmallBeck}, \(K\) is an equivalence—%
  in particular, \(\overline{\cat{A}_T}\) is semisimple.

  Assume now that \(\overline{\cat{A}_T}\) is semisimple.
  Then \(\overline{U_T}\) is exact by \cref{toomuchTV},
  and so it again satisfies the assumptions of \cref{thm:SmallBeck}.
  Thus, \(K\from \overline{\cat{A}_T} \iso \cat{A}^T\) is an equivalence;
  in particular, \(\cat{A}^T\) is semisimple.
\end{proof}

\begin{definition}\label{def:semisimple-monad}
  We say that a right exact monad \(T\from \cat{A} \to \cat{A}\) on an abelian category \(\cat{A}\) is \emph{semisimple}
  if it satisfies the equivalent conditions of \cref{prop:SemisimpleEM}.
\end{definition}

Since the opposite of an abelian category is abelian and the opposite of a semisimple category is semisimple,
an analogous result to \cref{prop:SemisimpleEM} holds for comonads,
and we define \emph{semisimple comonads} analogously.

\begin{proposition}\label{extendingsemisimple}
  Let \(T\from \cat{M} \to \cat{M}\) be a semisimple lax \(\cat{C}\)\hyp{}module monad on an abelian \(\cat{C}\)\hyp{}module category.
  Then \(T\) is extendable.
\end{proposition}
\begin{proof}
  By \cref{prop:SemisimpleEM}, the functor \(\overline{\iota}\from \overline{\cat{M}_T} \iso \cat{M}^T\) is an equivalence.
  Further, following \cref{cocompletingnonsense},
  \(\overline{\cat{M}_T}\) has a canonical structure of a \(\cat{C}\)\hyp{}module category,
  such that the inclusion \(\cat{M}_T \longhookrightarrow \overline{\cat{M}_T}\) is a strong \(\cat{C}\)\hyp{}module functor.
  Transporting the \(\cat{C}\)\hyp{}module structure along \(\overline{\iota}\) endows \(\iota\) with the structure of a strong \(\cat{C}\)\hyp{}module functor,
  proving the claim.
\end{proof}

Semisimplicity also ``transfers'' to the left adjoint of a monad.

\begin{lemma}\label{semisimpleleftandright}
  Let \(T\from \cat{A} \to \cat{A}\) be a right exact monad on an abelian category, and \(S\from \cat{A} \to \cat{A}\) a left adjoint to \(T\).
  Then \(T\) is semisimple if and only if \(S\) is semisimple.
  In that case, there is an equivalence \(\cat{A}^T \simeq \cat{A}^{S}\).
\end{lemma}
\begin{proof}
  By \cref{prop:SemisimpleEM}, \(T\) is semisimple if and only if so is \(\overline{\cat{A}_T}\).
  Moreover, the equivalence \(\cat{A}_T \simeq \cat{A}_S\) of \cref{einerKleinerKleisliÄquivalenz}
  extends to one on the Cauchy completions: \(\overline{\cat{A}_T} \simeq \overline{\cat{A}_S}\).
  Thus, \(\overline{\cat{A}_T}\) is semisimple if and only if \(\overline{\cat{A}_S}\) is so,
  which is the case if and only if \(S\) is semisimple.
  This establishes the first claim.

  The latter claim follows from the equivalences
  \[
    \cat{A}^T \,\simeq\, \overline{\cat{A}_T} \,\simeq\, \overline{\cat{A}_S} \,\simeq\, \cat{A}^{S}.\qedhere
  \]
\end{proof}

\begin{proposition}\label{ModuleEMSemisimple}
  Let \(T\from \cat{M} \to \cat{M}\) be a semisimple lax \(\cat{C}\)\hyp{}module monad on an abelian \(\cat{C}\)\hyp{}module category \(\cat{M}\),
  and let \(S\from \cat{M} \to \cat{M}\) be a left adjoint to \(T\).
  The equivalence \(\cat{M}^T \simeq \cat{M}^T\) of \cref{semisimpleleftandright} is a \(\cat{C}\)\hyp{}module equivalence,
  where \(\cat{M}^T\) and \(\cat{M}^{S}\) are endowed with the extended \(\cat{C}\)\hyp{}module structures
  of \(\cat{M}_T\) and \(\cat{M}_{S}\), respectively.
\end{proposition}
\begin{proof}
  Following the proof of \cref{extendingsemisimple},
  the functor \(\overline{\iota}\from \overline{\cat{M}_T} \iso \cat{M}^T\) is a \(\cat{C}\)\hyp{}module equivalence,
  and similarly for \(\cat{M}^{S}\) and \(\overline{\cat{M}_{S}}\).
  Further, from \cref{cocompletingnonsense} we obtain a \(\cat{C}\)\hyp{}module equivalence
  \(\overline{\cat{M}_T} \simeq \overline{\cat{M}_{S}}\)
  from the \(\cat{C}\)\hyp{}module equivalence \(\cat{M}_{T} \simeq \cat{M}_{S}\) of \cref{KleisliMonadCorrespondence}.

  The result follows by the following chain of \(\cat{C}\)\hyp{}module equivalences:
  \[
    \cat{M}^T \,\simeq\, \overline{\cat{M}_T} \,\simeq\, \overline{\cat{M}_{S}} \,\simeq\, \cat{M}^{S}.\qedhere
  \]
\end{proof}

\subsection{Linton coequalisers via multiactegories}\label{extendingeverything}

\textsc{There exists a quite general condition for extendability},
in which one can even explicitly write down the resulting \(\cat{C}\)\hyp{}module structure on the Eilenberg–Moore category.
Our aim now is to establish \cref{extendingalways,strongembeddingalways}.
Thus, for the rest of this section, we make the following assumptions.

\begin{hypothesis}\hfill
  \begin{itemize}[noitemsep]
    \item Let \(\cat{C}\) be an abelian monoidal and \(\cat{M}\) an abelian \(\cat{C}\)\hyp{}module category.
    \item The action functor \(\lact \from \cat{C} \kotimes \cat{M} \to \cat{M}\) is right exact in both variables.
    \item Let \(T\from \cat{M}\to \cat{M}\) be a right exact lax \(\cat{C}\)\hyp{}module monad.
  \end{itemize}
\end{hypothesis}

Note that, due to for example \cref{prop:abelian-EM-cat},
many of these assumptions are already necessary for the theory we would like to develop in this chapter.

In our presentation, we closely follow~\cite{aguiar18:monad},
where an analogous result for lax monoidal monads—rather than lax module monads—can be found.
As such, we shall often only describe the modifications necessary to make these results work in our case.

\begin{definition}
  The \emph{Linton coequaliser} \(x \blact m\) of \(x \in \cat{C}\) and \((m, \nabla_m) \in \cat{M}^T\) is:
  \[
    x \blact m \defeq \coeq
    \Big(\!\!
    \begin{mytikzcd}[ampersand replacement=\&]
      {T(x \lact Tm)} \& {T(x \lact m)}
      \arrow["{\raisebox{0.3em}{\(T(x \,\lact\, \nabla_{m})\)}}", shift left=2, from=1-1, to=1-2]
      \arrow["{\raisebox{-0.3em}{\(\mu_{x \lact m}\circ TT_{\mathsf{a};x,m}\)}}"', shift right=2, from=1-1, to=1-2]
    \end{mytikzcd}
    \!\!\Big).
  \]
\end{definition}

As \(\cat{M}^T\) is abelian due to \cref{prop:abelian-EM-cat},
by functoriality of colimits we obtain a functor \(\blank \blact \bblank\from \cat{C} \kotimes \cat{M}^T \to \cat{M}^T\).

\begin{lemma}\label{rightexactaction}
  For \(x \in \cat{C}\), the functor \(x \blact \blank\from \cat{M}^T \to \cat{M}^T\) is right exact.
\end{lemma}
\begin{proof}
  This immediately follows from the right exactness of \(T(x \lact T(\blank))\) and \(T(x \lact \blank)\),
  as well as the definition of the Linton coequaliser.
\end{proof}

\begin{lemma}\label{isounitality}
  There is a natural isomorphism \(1 \blact \blank \cong \Id_{\cat{M}^T}\).
\end{lemma}
\begin{proof}
  Recall that for all \(T\)\hyp{}algebras \(m \in \cat{M}^T\),
  \[
    \begin{mytikzcd}[ampersand replacement=\&]
      {T^{2}m} \& {Tm} \& m
      \arrow["{\mu_{m}}", shift left=1, from=1-1, to=1-2]
      \arrow["{T\nabla_m}"', shift right=1, from=1-1, to=1-2]
      \arrow["{\nabla_m}", from=1-2, to=1-3]
    \end{mytikzcd}
  \]
  is a coequaliser in \(\cat{M}^T\),
  created by the underlying split coequaliser in \(\cat{M}\).
  Now,
  \[
    \begin{mytikzcd}[ampersand replacement=\&]
      {T(1 \lact Tm)} \& {T(1 \lact m)} \\
      {T^{2}m} \& {Tm}
      \arrow["{T(1\lact \nabla_{m})}", shift left=1, from=1-1, to=1-2]
      \arrow["{\mu_{1\lact m} \circ T T_{\mathsf{a};1,m}}"', shift right=1, from=1-1, to=1-2]
      \arrow["{T \lambda}"', from=1-1, to=2-1]
      \arrow["\cong", from=1-1, to=2-1]
      \arrow["{T \lambda}", from=1-2, to=2-2]
      \arrow["\cong"', from=1-2, to=2-2]
      \arrow["{T \nabla_{m}}", shift left=1, from=2-1, to=2-2]
      \arrow["{\mu_{m}}"', shift right=1, from=2-1, to=2-2]
    \end{mytikzcd}
  \]
  defines an isomorphism in the category of parallel pairs of morphisms in \(\cat{M}^T\)%
  —the upper square commutes due to the naturality of the left unitor \(\lambda\) of \(\cat{M}\),
  and the lower by coherence for \(\blank \lact \bblank\).
  This isomorphism induces an isomorphism on the level of coequalisers, \(m \cong 1 \blact m\).
  Further, this construction of morphisms of parallel pairs is clearly natural in \(m\),
  establishing the naturality of the isomorphism of coequalisers.
\end{proof}

\begin{lemma}\label{actiononfree}
  For any \(m \in \cat{M}\), the coequaliser of
  \[
    \begin{mytikzcd}[ampersand replacement=\&]
      {T(x \lact T^{2}m)} \& {T(x\lact Tm)}
      \arrow["{T(x \lact \mu_{m})}", shift left=1, from=1-1, to=1-2]
      \arrow["{\mu_{x\lact Tm}\circ T T_{\mathsf{a};x,Tm}}"', shift right=1, from=1-1, to=1-2]
    \end{mytikzcd}
  \]
  in \(\cat{M}^T\) is natural in \(m\), split, and isomorphic to \(T(x \lact m)\).
  In other words the Linton coequaliser on a free \(T\)\hyp{}module is given by \(x \blact Tm \cong T(x \lact m)\).
\end{lemma}
\begin{proof}
  The splitting data is given by
  \[
    \begin{mytikzcd}[ampersand replacement=\&]
      {T(x \lact T^{2}m)} \& {T(x\lact Tm)} \& {T(x\lact m).}
      \arrow["{T(x \lact \mu_{m})}", from=1-1, to=1-2]
      \arrow["{\mu_{x\lact Tm}\circ T T_{\mathsf{a};x,Tm}}"', shift right=2, from=1-1, to=1-2]
      \arrow["{T(x \lact T\eta_{m})}"', curve={height=18pt}, from=1-2, to=1-1]
      \arrow["{\mu_{x \lact m}\circ T T_{\mathsf{a};x,m}}"', from=1-2, to=1-3]
      \arrow["{T(x \lact \eta_{m})}"', curve={height=18pt}, from=1-3, to=1-2]
    \end{mytikzcd}
    \qedhere
  \]
\end{proof}

\begin{lemma}\label{isoassociativity}
 For \(x,y \in \cat{C}\), there is a natural isomorphism \((x \otimes y) \blact \blank \cong x \blact (y \blact \blank)\).
\end{lemma}
\begin{proof}
  We calculate
  \[
    \begin{aligned}
      x \blact (y \blact m)
      &= x \blact \coeq(T(y \lact \nabla_{m}),\, \mu_{x \lact m} \circ  T T_{\mathsf{a};y,m})\\
      &\cong \coeq(x \blact T(y \lact \nabla_{m}),\, x \blact \mu_{x \lact m} \circ  T T_{\mathsf{a};y,m})\\
      &\cong \coeq(T(x\lact (y \lact \nabla_{m})),\, \mu_{y \lact x \lact m}\circ T(T_{\mathsf{a}; x, y\lact m} \circ y \lact T_{\mathsf{a};x,m}))\\
      &\cong \coeq(T((x \otimes y) \lact \nabla_{m}),\, \mu_{(x \otimes y)\lact Tm} \circ T T_{\mathsf{a}; x \otimes y, Tm}) \\
      &= (x \otimes y) \blact m,
    \end{aligned}
  \]
  where the first isomorphism follows from \cref{rightexactaction},
  the second one is due to \cref{actiononfree},
  and the third one follows by coherence of \(\blank \lact \bblank\).
\end{proof}

In view of \cref{isounitality,isoassociativity},
all that remains in order to establish that \(\blank \blact \bblank\) defines a \(\cat{C}\)\hyp{}module category structure on \(\cat{M}^T\) is the verification of coherence axioms.
While this can be accomplished by diagram chasing—see~\cite[Theorem~2.6.4]{seal13:tensor}—we choose to give a more formal argument,
mimicking the multicategorical approach used by~\cite{aguiar18:monad} in the analogous case of monoidal monads.
More specifically, we shall prove the following two results in the remainder of this section.

\begin{theorem}\label{extendingalways}
 The functor \(\blank \blact \bblank \from \cat{C} \kotimes \cat{M}^T \to \cat{M}^T\) defines a left \(\cat{C}\)\hyp{}module category structure on \(\cat{M}^T\).
\end{theorem}

Assuming \cref{extendingalways},
we have \(x \blact Tm = T(x \lact m)\),
and so the image of the inclusion \(\iota\from \cat{M}_T \to \cat{M}^T\) of \cref{denotedbyiota}
is a \(\cat{C}\)\hyp{}module subcategory,
where \(\cat{M}_T\) is endowed with the action given in the proof of \cref{laxKlCorrespondence}.

\begin{theorem}\label{strongembeddingalways}
 The unique embedding \(\iota \from (\cat{M}_T, \lact_{\cat{M}_T}) \to (\cat{M}^T, \blact)\) can be equipped with the structure of a strong \(\cat{C}\)\hyp{}module functor.
\end{theorem}

Let us now recall the theory of multicategories,
see~\cite{hermida00:repres,leinster04:higher}.

\begin{definition}\label{def:multicategory}
  A (locally small) \emph{multicategory} \(\mathbf{C}\) consists of
  \begin{itemize}
    \item a class \(\Ob\mathbf{C}\) of objects of \(\mathbf{C}\), where we write \(x \in \mathbf{C}\) for \(x \in \Ob\mathbf{C}\);
    \item for any finite sequence \((x_{1},\ldots, x_{n}, y)\) of objects in \(\mathbf{C}\),
    a set of \emph{multimorphisms} \(\mathbf{C}(x_{1},\ldots, x_{n}; y)\) from \((x_{1},\ldots,x_{n})\) to \(y\);
    \item for every \(x \in \mathbf{C}\), an identity multimorphism \(\id_{x} \in \mathbf{C}(x;x)\); and
    \item for any \(y \in \mathrm{Ob}(\mathbf{C})\), \({(x_{i})}_{i=1}^{n} \in {\mathrm{Ob}(\mathbf{C})}^{n}\),
    as well as
    \[
      \big((u_1^1, \dots, u_1^{k_1}), \dots, (u_n^1, \dots, u_n^{k_n})\big) \in {\mathrm{Ob}(\mathbf{C})}^{k_{1} + \dots + k_n},
    \]
    a composition operation
    \[
      \begin{aligned}
        \mathbf{C}&(x_{1},\ldots,x_{n}; y) \times \mathbf{C}(u_{1}^{1},\ldots, u_{1}^{k_{1}}; x_{1}) \times \cdots \times \mathbf{C}(u_{n}^{1},\ldots, u_{n}^{k_{n}}; x_{n}) \\
         &\to \mathbf{C}(u_{1}^{1},\ldots,u_{n}^{k_{n}}; y),
      \end{aligned}
    \]
  \end{itemize}
  subject to natural associativity and unitality conditions; see for example~\cite[Section~2.1]{leinster04:higher} or~\cite[Definition~2.1]{hermida00:repres}.
\end{definition}

Since we want our multicategories and multiactegories to correspond to \(\Bbbk\)\hyp{}linear monoidal and module categories,
we should replace sets of multimorphisms with vector spaces,
and the Cartesian products with tensor products.
However, since our aim is to show \cref{extendingalways},
which can be verified on the level of the underlying (ordinary) categories,
this is not essential.
The following definition is a non-skew special case of~\cite[Definition~6.9]{arkor24}.

\begin{definition}\label{def:multiactegory}
  Let \(\mathbf{C}\) be a (locally small) multicategory.
  A (locally small) \emph{left multiactegory} \(\mathbf{M}\) over \(\mathbf{C}\) has as data:
  \begin{itemize}
    \item a class \(\Ob\mathbf{M}\) of objects of \(\mathbf{M}\),
    where we write \(m \in \mathbf{M}\) for \(m \in \mathrm{Ob}\mathbf{M}\);
    \item for any finite (possibly empty) sequence \((x_{1},\ldots, x_{n})\) of objects in \(\mathbf{C}\)
    and a pair of objects \((m, m')\) of \(\mathbf{M}\),
    a set \(\mathbf{M}(x_{1},\ldots, x_{n}; m; m')\) of \emph{multimorphisms} from \((x_{1},\ldots,x_{n}; m)\) to \(m'\);
    \item for every \(m \in \mathbf{M}\), an identity multimorphism \(\mathrm{id}_{m} \in \mathbf{M}(m;m)\); and
    \item for any \(x \in \mathbf{C}\), \(m, m', \ell \in \mathbf{M}\),
    \({(x_{i})}_{i=1}^{n} \in {\mathrm{Ob}(\mathbf{C})}^{n}\),
    and
    \[
      \big((u_1^1, \dots, u_1^{k_1}), \dots, (u_n^1, \dots, u_n^{k_n})\big) \in {\mathrm{Ob}(\mathbf{C})}^{k_{1} + \dots + k_n},
    \]
    a composition operation
    \begin{equation}\label{eq:multiactcompose}
      \mathscale{0.95}{\begin{aligned}
        \mathbf{M}&(x_{1},\ldots,x_{n}; m; m')
                    \times \mathbf{C}(u_{1}^{1},\ldots, u_{1}^{k_{1}};x_{1})
                    \times \cdots
                    \times \mathbf{C}(u_{n}^{1},\ldots, u_{n}^{k_{n}};x_{n}) \\
                  &\times \mathbf{M}(u_{n+1}^{1},\ldots,u_{n+1}^{k_{n+1}}; \ell; m)
                    \to \mathbf{M}(u_{1}^{1},\ldots,u_{n+1}^{k_{n+1}}; \ell; m').
      \end{aligned}}
    \end{equation}
  \end{itemize}
  This data has to satisfy the (non-marked) associativity and unitality axioms of~\cite[Definitions~4.4 and~6.9]{arkor24}.
\end{definition}

\begin{notation}
  From now on, we will often abbreviate sequences of objects using the vector notation similar to~\cite{hermida00:repres};
  \eg, writing \(\vec{x}\) for \((x_{1},\ldots, x_{n})\)
  and \((\vec{x_{1}}, \vec{x_{2}})\) for \((x_{1}^{1},\ldots, x_{1}^{n_{1}}, x_{2}^{1}, \ldots, x_{2}^{n_{2}})\).
  We will use this notation for clarity and brevity, when keeping track of the length of the tuples is not required.
  We also fix the notation for the remainder of this section,
  letting \(\mathbf{C}\) be a multicategory and letting \(\mathbf{M}\) be a \(\mathbf{C}\)\hyp{}multiactegory.
\end{notation}

Recall from~\cite[Chapter~8]{hermida00:repres} that a \emph{pre-universal arrow} for a tuple \(\vec{x}\) of objects of \(\mathbf{C}\)
consists of an object \(\otimes(\vec{x})\) of \(\mathbf{C}\)
together with a multimorphism \(\pi \in \mathbf{C}(\vec{x};\otimes(\vec{x}))\),
inducing isomorphisms \(\mathbf{C}(\otimes(\vec{x});y) \iso \mathbf{C}(\vec{x};y)\).
If \(\pi\) further induces isomorphisms
\[
  \mathbf{C}(\vec{z}, \otimes(\vec{x}), \vec{z'}; y) \iso \mathbf{C}(\vec{z}, \vec{x}, \vec{z'}; y),
\]
we say that \(\pi\) is \emph{universal}.
If any sequence of objects of \(\mathbf{C}\) is the domain of a universal arrow,
we say that \(\mathbf{C}\) is \emph{representable}.
A \emph{representation} of \(\mathbf{C}\) is a choice of a universal arrow from every sequence of objects in \(\mathbf{C}\).

\begin{definition}
  Let \(\mathbf{M}\) be a \(\mathbf{C}\)\hyp{}multiactegory.
  A \emph{pre-universal arrow} for a tuple \(\vec{x}\) of objects of \(\mathbf{C}\)
  and \(m \in \mathbf{M}\)
  consists of an object \(\lact(\vec{x}; m)\),
  together with a multimorphism \(\pi \in \mathbf{M}(\vec{x};m; \lact(\vec{x};m))\),
  inducing isomorphisms
  \[
    \mathbf{M}(\lact(\vec{x};m);m') \iso \mathbf{M}(\vec{x};m;m').
  \]

  If \(\pi\) further induces isomorphisms
  \[
    \mathbf{M}(\vec{y},\lact(\vec{x};m);m') \iso \mathbf{M}(\vec{y},\vec{x};m;m'),
  \]
  then \(\pi\) is said to be \emph{universal}.
\end{definition}

\begin{definition}
  We say that a multimorphism \(\rho \in \mathbf{C}(\vec{x}; y)\) is \emph{pre-universal in \(\mathbf{M}\)}
  if it induces isomorphisms \(\mathbf{M}(y; m; n) \iso \mathbf{M}(\vec{x}; m; n)\).
  It is \emph{universal in \(\mathbf{M}\)}, if it also induces isomorphisms
  \(\mathbf{M}(\vec{z}, \vec{x}, \vec{w}; m; n) \iso \mathbf{M}(\vec{z}, y, \vec{w}; m; n)\).
\end{definition}

\begin{definition}\label{representablemultiactegory}
  For \(\mathbf{C}\) representable and a representation \(\mathsf{R}\) of \(\mathbf{C}\),
  we say that \(\mathbf{M}\) is a \emph{representable with respect to \(\mathsf{R}\)},
  if for every sequence \(\vec{x}\) of objects of \(\mathbf{C}\) and every object \(m \in \mathbf{M}\),
  there is a universal arrow \(\pi_{(\vec{x};m)} \in \mathbf{M}(\vec{x};m; \lact(\vec{x},m))\),
  and all morphisms of \(\mathsf{R}\) are universal in \(\mathbf{M}\).
\end{definition}

Recall from~\cite[Proposition~8.5]{hermida00:repres} that universal arrows in a multicategory \(\mathbf{C}\) are closed under composition.
Further, if every sequence of objects in \(\mathbf{C}\) is the domain of a pre-universal arrow and pre-universal arrows are closed under composition,
then a pre-universal arrow in \(\mathbf{C}\) is universal.
A similar result holds for multiactegories.

\begin{lemma}\hfill
  \begin{enumerate}[noitemsep]
    \item Composition of universal arrows of \(\mathbf{C}\) in \(\mathbf{M}\)
    and universal arrows of \(\mathbf{M}\) yields universal arrows in \(\mathbf{M}\).
    \item If composition of pre-universal arrows of \(\mathbf{C}\) in \(\mathbf{M}\)
    and pre-universal arrows of \(\mathbf{M}\) yields pre-universal arrows in \(\mathbf{M}\),
    then a pre-universal morphism in \(\mathbf{M}\) is universal.
    \item If pre-universal arrows in \(\mathbf{C}\) act pre-universally in \(\mathbf{M}\),
    and pre-universal arrows of \(\mathbf{C}\) and \(\mathbf{M}\) are closed under composition in \(\mathbf{M}\),
    then \(\mathbf{M}\) is representable.
  \end{enumerate}
\end{lemma}

\begin{proof}
  For the first part, let \((\pi_{0}, \rho_{1},\ldots, \rho_{n}, \pi_{n+1})\) be a sequence composable in \(\mathbf{M}\),
  with \(\pi_{0},\pi_{n+1}\) in \(\mathbf{M}\) and \(\rho_{1},\ldots, \rho_{n} \in \mathbf{C}\),
  as indicated in \cref{eq:multiactcompose},
  all of whose elements are universal in \(\mathbf{M}\).
  The composite \(\pi_{0} \circ (\rho_{1},\ldots, \rho_{n}, \pi_{n+1})\) gives rise to the sequence of isomorphisms
  \[
    \begin{aligned}
      \mathbf{M}&(\vec{x}, \lact(\otimes(\vec{y}_{1}),\ldots, \otimes(\vec{y}_{n});\lact(\vec{y}_{n+1},m)),\blank) \\
                  &\cong \mathbf{M}(\vec{x}, \otimes(\vec{y}_{1}), \ldots, \otimes(\vec{y}_{n});\lact(\vec{y}_{n+1},m);\blank)
                \cong \mathbf{M}(\vec{x},\ldots, \vec{y}_{1},\ldots,\vec{y}_{n+1};m;\blank),
    \end{aligned}
  \]
  proving the universality of the composite.

  For the second part, consider the following chain of maps:
  \begin{align*}
    \mathbf{M}(\vec{x},\vec{y},\vec{z};m;\blank)
    & \rightarrow \mathbf{M}(\vec{x},\otimes(\vec{y}),\vec{z};m;\blank)
      \rightarrow \mathbf{M}(\otimes(\vec{x}),\otimes(\vec{y}),\otimes(\vec{z});m;\blank) \\
    & \rightarrow \mathbf{M}(\lact(\otimes(\vec{x}),\otimes(\vec{y}),\otimes(\vec{z});m);\blank),
  \end{align*}
  where the first morphism is induced by the pre-universal arrow \(\rho_{\vec{x}}\) of \(\mathbf{C}\),
  the second one by the pre-universal arrows \(\rho_{\vec{y}}\) and \(\rho_{\vec{z}}\) of \(\mathbf{C}\),
  and the third is induced by the universal arrow \(\pi_{(\vec{y},\vec{x},\vec{z}; m)}\).
  The composite of all three maps is an isomorphism%
  —being induced by a composite of pre-universal arrows—%
  and hence a pre-universal arrow.
  The same holds for the composite of the latter two maps.
  These two observations establish the invertibility of the first map.

  The third part is established by applying the proof of the second part above to the case of pre-universal arrows of \(\mathbf{C}\).
\end{proof}

Assume \(\mathbf{C}\) to be representable, let \(\mathsf{R}\) be a representation of \(\mathbf{C}\),
and suppose that \(\mathbf{M}\) is representable with respect to \(\mathsf{R}\).
Given a sequence \(x_{1},\ldots,x_{n}\) of objects in \(\mathbf{C}\) and an object \(m\) in \(\mathbf{M}\),
the set of ways in which we may compose the universal arrows in the given representations to a universal arrow with domain \((x_{1},\ldots,x_{n};m)\)
is canonically in bijection with the set of parenthesisations of the word \(x_{1}\cdots x_{n}m\)
into subwords of length at least two.
Let \(\diamond\) and \(\diamond'\) be two such parenthesisations.
We denote the codomains of the corresponding universal arrows by \(\diamond(x_{1},\ldots,x_{n};m)\) and \(\diamond'(x_{1},\ldots,x_{n};m)\),
and the universal arrows themselves by \(\pi_{\diamond(x_{1},\ldots,x_{n};m)}\) and \(\pi_{\diamond'(x_{1},\ldots,x_{n};m)}\).
The universality of \(\pi_{\diamond(x_{1},\ldots,x_{n};m)}\) entails the existence of
a unique arrow
\[
  \alpha_{\vec{x},m}^{\diamond,\diamond'}\from \diamond(x_{1},\ldots,x_{n};m)\to \diamond'(x_{1},\ldots,x_{n};m),
\]
satisfying \(\alpha_{\vec{x},m}^{\diamond,\diamond'} \circ \pi_{\diamond(\vec{x},m)} = \pi_{\diamond'(\vec{x},m)}\),
which is invertible with inverse \(\alpha_{\vec{x},m}^{\diamond',\diamond}\).

Further, \(\alpha_{\vec{x},m}^{\diamond,\diamond''} = \alpha_{\vec{x},m}^{\diamond',\diamond''} \circ \alpha_{\vec{x},m}^{\diamond,\diamond'}\).
Thus, for any sequence \(x_{1},\ldots,x_{n},m\) we obtain
an indiscrete category\note{%
  A category is called \emph{indiscrete} if there exists a unique morphism from every object to every other object.%
}
whose objects are parenthesisations of \(x_{1}\cdots x_{n}m\)
and \(\mathrm{Hom}(\diamond,\diamond') = \{\alpha_{\vec{x},m}^{\diamond,\diamond'}\}\).
The arrow \(\alpha_{\vec{x},m}^{\diamond,\diamond'}\) is natural in \(m\) and \(x_{i}\) for all \(i\).

\begin{remark}\label{rmk:monoidal-from-multi}
  Recall that given a representable multicategory \(\mathbf{C}\) with a fixed representation,
  there is a monoidal category \(\cat{C}\) defined by \(\Ob\cat{C} \defeq \Ob\mathbf{C}\),
  \(\cat{C}(x,y) \defeq \mathbf{C}(x;y)\),
  and \(x \otimes y \defeq \otimes(x,y)\),
  the codomain of the universal arrow \(\pi_{\otimes(x,y)}\) from \((x,y)\) in the representation of \(\mathbf{C}\).
  The associator of \(\cat{C}\) is defined by the
  \(\alpha_{w,x,y}^{(wx)y,w(xy)}\) described above;
  one obtains similar arrows for unitality.

  The indiscreteness of the category of morphisms \(\alpha^{\diamond,\diamond'}\)
  associated to the parenthesisations of words in \(\mathbf{C}\)
  implies the commutativity of the pentagon and triangle diagrams,
  establishing the coherence and well-definedness of \(\cat{C}\).
\end{remark}

The indiscrete category described above establishes
the well-definedness of the \(\cat{C}\)\hyp{}module category \(\cat{M}\) of the following definition.

\begin{definition}\label{associatedrepresented}
  Let \(\mathbf{M}\) be a representable multiactegory over a representable multicategory \(\mathbf{C}\),
  with fixed representations for both.
  Let \(\cat{C}\) be the monoidal category associated to the representation of \(\mathbf{C}\).
  The \emph{\(\cat{C}\)\hyp{}module category \(\cat{M}\) associated to the fixed representation of \(\mathbf{M}\)} is defined by setting \(\Ob\mathbf{M} \defeq \Ob\cat{M}\),
  \(\cat{M}(m,n) \defeq \mathbf{M}(m; n)\),
  and \(x \lact m \defeq \lact(x,m)\), the codomain of the universal arrow \(\pi_{\lact(x,m)}\).
  The associator \((x \otimes y) \lact m \iso x \lact (y \lact m)\) is given by the arrow \(\alpha_{x,y,m}^{(xy)m, x(ym)}\),
  and similarly for the unitors.
\end{definition}

Having established sufficient analogues of~\cite{hermida00:repres,leinster04:higher},
we continue with analogues of~\cite[Section~3]{aguiar18:monad}.
For \(x_i \in \cat{C}\) and \(m \in \cat{M}\), write
\[
  \varphi \from x_{1}\lact \cdots \lact x_{n} \lact Tm \to T(x_{1}\lact \cdots \lact x_{n} \lact m)
\]
for the morphism obtained by repeated application of the lax structure of \(T\),
where \(x_{1}\lact \cdots \lact x_{n} \lact m\) is to be read as \(x_{1} \lact (\cdots \lact (x_{n} \lact m))\).

We will also consider the coequaliser diagram
\begin{equation}\label[diagram]{unicharacterize}
  \mathscale{0.98}{\begin{mytikzcd}[ampersand replacement=\&]
    {T(x_{1} \lact \cdots \lact x_{n} \lact Tm)} \& {T^{2}(x_{1} \lact \cdots \lact x_{n} \lact m)} \& {T(x_{1} \lact \cdots \lact x_{n} \lact m)}
    \arrow["{T\varphi}"', from=1-1, to=1-2]
    \arrow["{T(x_{1}\lact\cdots \lact x_{n}\lact \nabla_{m})}", curve={height=-18pt}, from=1-1, to=1-3]
    \arrow["\mu"', from=1-2, to=1-3]
  \end{mytikzcd}}
\end{equation}
which is reflexive, with
\(T(x_{1} \lact \cdots \lact x_{n} \lact m) \xrightarrow{T(x_{1} \lact \cdots \lact x_{n} \lact \eta_{m})} T(x_{1} \lact \cdots \lact x_{n} \lact Tm)\)
as a common section.

\begin{definition}
  A morphism \(f\from x_{1}\lact \cdots \lact x_{n} \lact m \to m'\)
  is called \emph{\(n\)\hyp{}multilinear}
  if the following diagram commutes:
  \[
    \begin{mytikzcd}[ampersand replacement=\&]
      {x_{1}\lact \cdots \lact x_{n} \lact Tm} \&\& {T(x_{1}\lact \cdots \lact x_{n}\lact m)} \\
      \&\& {Tm'} \\
      {x_{1}\lact \cdots \lact x_{n} \lact m} \&\& m'
      \arrow["\varphi", from=1-1, to=1-3]
      \arrow["{x_{1}\lact \cdots \lact x_{n} \lact \nabla_{m}}"', from=1-1, to=3-1]
      \arrow["{Tf}", from=1-3, to=2-3]
      \arrow["{\nabla_{m'}}", from=2-3, to=3-3]
      \arrow["f"', from=3-1, to=3-3]
    \end{mytikzcd}
  \]
\end{definition}

\begin{remark}
  For a monoidal category \(\cat{C}\), there exists a representable multicategory \(\mathbf{C}\),
  with \(\mathbf{C}(x_{1},\ldots,x_{n};y) \defeq \cat{C}(x_{1}\otimes (\cdots (x_{n-1}\otimes x_{n})\ldots),y)\).
  Then,
  a multimorphism \(\pi \in \mathbf{C}(x_{1},\ldots,x_{n}; y)\) is universal if and only if
  it is an isomorphism \(\pi\from x_{1} \otimes \cdots \otimes x_{n} \iso y\).

  Similarly, given a \(\cat{C}\)\hyp{}module category,
  one obtains a representable multiactegory \(\mathbf{M}\) over \(\mathbf{C}\)
  satisfying \(\mathbf{M}(x_{1},\ldots,x_{n}; m; n) = \cat{M}(x_{1}\lact \cdots \lact x_{n} \lact m, m')\).
  A universal arrow \(\pi \in \mathbf{M}(x_{1},\ldots,x_{n}; m; m')\) is an isomorphism \(\pi\from x_{1} \lact \cdots \lact m \iso m'\).
\end{remark}

\begin{lemma}
  For a \(\cat{C}\)\hyp{}module category \(\cat{M}\),
  the collection of multilinear morphisms of \(\mathbf{M}\)
  is stable under multicomposition with morphisms of \(\mathbf{C}\),
  forming a \(\cat{C}\)\hyp{}multiactegory that we denote by \(\mathbf{M}^{\varphi}\).
\end{lemma}
\begin{proof}
  For \(i = 1,\ldots, n\), consider sequences \(\vec{z_{i}}\);
  morphisms \(\otimes(z_{i}) \xrightarrow{g_{i}} x_{i}\),
  where \(\otimes(z_{i})\) again is given the rightmost parenthesisation;
  and a multilinear morphism \(f\from x_1 \lact \dots \lact x_n \lact m \to m'\) of \(\mathbf{M}\).
  The claim follows immediately from the commutativity of
  \[
    \begin{mytikzcd}[ampersand replacement=\&]
      {\otimes(\vec{z_{1}})\lact \cdots \lact \otimes(\vec{z_{n}}) \lact Tm} \&\& {T(\otimes(\vec{z_{1}})\lact \cdots \lact \otimes(\vec{z_{n}})\lact m)} \\
      {x_{1}\lact \cdots \lact x_{n} \lact Tm} \&\& {T(x_{1}\lact \cdots \lact x_{n}\lact m)} \\
      \&\& {Tm'} \\
      {x_{1}\lact \cdots \lact x_{n} \lact m} \&\& m'
      \arrow["\varphi", from=1-1, to=1-3]
      \arrow["{g_{1}\lact\cdots \lact g_{n} \lact Tm}", from=1-1, to=2-1]
      \arrow["{T(g_{1}\lact\cdots \lact g_{n} \lact m)}", from=1-3, to=2-3]
      \arrow["\varphi", from=2-1, to=2-3]
      \arrow["{x_{1}\lact \cdots \lact x_{n} \lact \nabla_m}", from=2-1, to=4-1]
      \arrow["{Tf}"', from=2-3, to=3-3]
      \arrow["{\nabla_{m'}}"', from=3-3, to=4-3]
      \arrow["f", from=4-1, to=4-3]
    \end{mytikzcd}
  \]
  where the top face commutes by naturality of the action of \(\cat{C}\) in \(\cat{M}\).
\end{proof}

Note that \(\mathbf{M}^{\varphi}\) being representable would prove \cref{extendingalways},
since then the \(\cat{C}\)\hyp{}module category can be equipped with coherent associators;
see \cref{rmk:monoidal-from-multi,associatedrepresented}.

\begin{lemma}\label{multiuni}
  A map \(f\) is \(n\)\hyp{}multilinear if and only if
  the map
  \[
    \underline{f} \from T(x_1 \lact \dots \lact x_n \lact m) \xrightarrow{\ Tf\ } Tm' \xrightarrow{\ \nabla_{m'}\ } m'
  \]
  coequalises \cref{unicharacterize},
  and \(f\) is pre-universal in \(\mathbf{M}^{\varphi}\) if and only if
  \(\underline{f}\) is the (universal) coequaliser of \emph{ibid}.
\end{lemma}
\begin{proof}
  Consider the following diagram:
  \[
    \begin{mytikzcd}[ampersand replacement=\&]
      {T(x_{1}\lact \cdots \lact x_{n}\lact Tm)} \&\& {T^{2}(x_{1}\lact \cdots \lact x_{n}\lact m)} \\
      \&\& {T^{2}m'} \\
      {T(x_{1}\lact \cdots \lact x_{n}\lact m)} \&\& {Tm'} \& {Tm'} \\
      \&\& m'
      \arrow["{T\varphi}", from=1-1, to=1-3]
      \arrow["{T(x_{1}\lact \cdots\lact \nabla_{m})}"', from=1-1, to=3-1]
      \arrow["{T^{2}f}", from=1-3, to=2-3]
      \arrow[""{name=0, anchor=center, inner sep=0}, "\mu", from=1-3, to=3-1]
      \arrow["{\mu_{m'}}", from=2-3, to=3-3]
      \arrow[""{name=1, anchor=center, inner sep=0}, "{T(\nabla_{m'})}", from=2-3, to=3-4]
      \arrow["{Tf}", from=3-1, to=3-3]
      \arrow["{\underline{f}}"', from=3-1, to=4-3]
      \arrow["{\nabla_{m'}}", from=3-3, to=4-3]
      \arrow["{\nabla_{m'}}", from=3-4, to=4-3]
      \arrow["{(1)}"{description}, draw=none, from=1-1, to=0]
      \arrow["{(2)}"{description}, shift left=5, draw=none, from=0, to=3-3]
      \arrow["{(3)}"{description, pos=0.4}, draw=none, from=1, to=4-3]
    \end{mytikzcd}
  \]
  Coequalising the pair~\eqref{unicharacterize} is precisely coequalising face \((1)\) in the above diagram.
  Thus, since face \((2)\) commutes, it is equivalent to \(\nabla_{m'}\) coequalising the square formed jointly by faces \((1)\) and \((2)\).

  Observe that, using commutativity of face \((3)\),
  this square postcomposed by \(\nabla_{m'}\) gives precisely the pair of morphisms in \(\cat{M}^T(T(x_{1}\lact \cdots \lact Tm), m')\)
  corresponding to the pair of morphisms in \(\cat{M}(x_{1}\lact \cdots \lact T(m), m')\),
  defining multilinearity of \(f\),
  under \(\cat{M}^T(T(\blank),\bblank) \cong \cat{M}(\blank,\bblank)\).
  Thus, both conditions in the first statement are equivalent to the commutativity of the diagram.

  The second statement is shown analogously to~\cite[Lemma~3.5]{aguiar18:monad}
\end{proof}

Finally, in order to establish \cref{extendingalways,strongembeddingalways},
we need an analogue of~\cite[Proposition~3.13]{aguiar18:monad}.

\begin{proposition}\label{hillbility}
 The multicategory \(\mathbf{M}^{\varphi}\) is representable.
\end{proposition}

\begin{proof}
  Let \((x_{1},\ldots,x_{n},m')\) be a sequence of objects, with \(x_{i} \in \cat{C}\) and \(m' \in \cat{M}\).
  The coequalizer of Diagram~\eqref{unicharacterize} defines a pre-universal map in \(\mathbf{M}^{\varphi}\) with domain \((x_{1},\ldots,x_{n};m')\), by Lemma~\ref{multiuni}.
  It remains to show that a composition of pre-universal maps is pre-universal.
  Thus, consider sequences \(\vec{z_{1}},\ldots,\vec{z_{n}}\) and \(\vec{y}\), universal multimorphisms \(h_{i}\from \vec{z_{i}} \to x_{i}\) of \(\mathbf{C}\), and universal multimorphisms \(f\from (\vec{y},m) \to m'\) and \(g\from (\vec{x},m') \to m''\).

  Behold the following diagram:
  \[
    \mathscale{0.8}{\begin{mytikzcd}[ampersand replacement=\&]
        \&\& {\vec{z_{1}} \lact \cdots \vec{z_{n}}\lact T^{2}(\vec{y} \lact m)} \\
        {\vec{z_{1}} \lact \cdots \vec{z_{n}}\lact T(\vec{y} \lact T(m))} \&\&\&\& {\vec{z_{1}} \lact \cdots \vec{z_{n}}\lact T(\vec{y} \lact m)} \\
        \&\& {T^{2}(\vec{z_{1}} \lact \cdots \vec{z_{n}}\lact \vec{y} \lact m)} \\
        {T(\vec{z_{1}} \lact \cdots \vec{z_{n}}\lact \vec{y} \lact T(m))} \&\&\&\& {T(\vec{z_{1}} \lact \cdots \vec{z_{n}}\lact \vec{y} \lact m)}
        \arrow["{z_{1} \lact \cdots z_{n}\lact \mu_{y\lact m}}"{pos=0.6}, from=1-3, to=2-5]
        \arrow["{\vec{z_{1}} \lact \cdots \vec{z_{n}}\lact T(\varphi)}"{pos=0.4}, from=2-1, to=1-3]
        \arrow["{\vec{z_{1}} \lact \cdots \vec{z_{n}}\lact T(\vec{y} \lact \nabla_{m})}", from=2-1, to=2-5]
        \arrow["\varphi"', from=2-1, to=4-1]
        \arrow["\varphi"', from=2-5, to=4-5]
        \arrow["{\mu_{z_{1}\lact \cdots \lact m}}", from=3-3, to=4-5]
        \arrow["{T\varphi}", from=4-1, to=3-3]
        \arrow["{T(\vec{z_{1}} \lact \cdots \vec{z_{n}}\lact \vec{y} \lact \nabla_{m})}"', from=4-1, to=4-5]
      \end{mytikzcd}}
  \]
  Similarly to the proof of~\cite[Proposition~3.13]{aguiar18:monad},
  this is a morphism from the top triangle-shaped diagram to the bottom one,
  where the top is the action of \((\vec{z_{1}},\ldots,\vec{z_{n}})\)
  on the Linton pair for \((\vec{y},m)\),
  while the bottom is the Linton pair for \((\vec{z_{1}},\ldots,\vec{z_{n}},\vec{y};m)\).

  Now, let \(X \in \mathbf{M}^{\varphi}(\vec{z_{1}},\ldots,\vec{z_{n}},\vec{y};m;\ell)\).
  Consider the diagram
  \[
    \mathscale{0.98}{\begin{mytikzcd}[ampersand replacement=\&]
      {\vec{z_{1}} \lact \cdots \vec{z_{n}}\lact T(\vec{y} \lact m)} \&\& {\vec{z_{1}} \lact \cdots \vec{z_{n}}\lact m'} \&\& {x_{1} \lact \cdots x_{n}\lact m'} \\
      \\
      {T(\vec{z_{1}} \lact \cdots \vec{z_{n}}\lact \vec{y} \lact m)} \&\& {\ell} \&\& m''
      \arrow["{z_{1}\lact\cdots \lact \underline{f}}", from=1-1, to=1-3]
      \arrow["\varphi"', from=1-1, to=3-1]
      \arrow["{h_{1} \lact \cdots h_{n}\lact m'}", from=1-3, to=1-5]
      \arrow["\simeq"', from=1-3, to=1-5]
      \arrow["{\exists! \widetilde{X}}", dashed, from=1-3, to=3-3]
      \arrow["g", from=1-5, to=3-5]
      \arrow["{\underline{X}}", from=3-1, to=3-3]
      \arrow["{\exists! \widehat{X}}", dotted, from=3-5, to=3-3]
    \end{mytikzcd}}
  \]
  where \(\widetilde{X}\) exists since, by \cref{multiuni},
  \(\underline{X}\) coequalises the bottom Linton pair in the penultimate diagram above,
  and \({h_{1}\lact\cdots \lact \underline{f}}\) is the coequaliser of the top Linton pair in the same diagram.

  Similarly to~\cite[Lemma~3.11]{aguiar18:monad} one shows that \(\widetilde{X}\) is multilinear,
  and thus so is \(h_{1}^{-1} \lact \cdots \lact h_{n}^{-1} \lact m'\),
  since \(h_{i}\) is invertible for all \(i\), being a pre-universal, and hence universal, arrow of \(\mathbf{C}\).
  This proves the existence of \(\widehat{X}\).

  Thus, the multilinear morphism \(X\) factors through \(g \circ (h_{1},\ldots, h_{n},f)\),
  and this factorisation is unique by uniqueness of \(\widetilde{X}\) and of \(\widehat{X}\).
  This establishes the pre-universality of the composition \(g \circ (h_{1},\ldots, h_{n}, f)\).
\end{proof}

\section{Internal projective and injective objects}\label{sec:internal-projectives-and-injectives}

\begin{definition}\label{def:c-projective}
  Let \(\cat{C}\) be an abelian monoidal category and let \(\cat{M}\) be an abelian \(\cat{C}\)\hyp{}module category.
  An object \(M \in \cat{M}\) is said to be \emph{\(\cat{C}\)\hyp{}projective} if,
  for any projective object \(P \in \cat{C}\),
  the object \(P \lact M \in \cat{M}\) is projective.
  Similarly, \(M\) is \emph{\(\cat{C}\)\hyp{}injective} if,
  for any injective object \(I \in \cat{C}\),
  the object \(I \lact M \in \cat{M}\) is injective.
\end{definition}

The next result connects
the formal definition of a \(\cat{C}\)\hyp{}projective
with the ordinary notion of projective object in an abelian category.

\begin{proposition}\label{Cprojectives}
  Assume that \(\cat{C}\) and \(\cat{M}\) have enough projectives.
  Let \(M\) be a closed object of \(\cat{M}\).
  Then \(\hom{M,\blank}\) is right exact if and only if \(M\) is \(\cat{C}\)\hyp{}projective.
\end{proposition}
\begin{proof}
  If \(\hom{M, \blank}\) is right exact and \(P\) is a projective object of \(\cat{C}\).
  Then \(P \lact M\) is projective because
  \(\cat{M}(P \lact M, \blank) \cong \cat{C}(P,\hom{M, \blank})\),
  which is a composite of right exact functors, hence itself a right exact functor.

  Assume that \(M\) is \(\cat{C}\)\hyp{}projective.
  Let \(\colim_{j} N_{j}\) be a finite colimit in \(\cat{M}\),
  and \(X \in \cat{C}\).
  Since \(\cat{C}\) has enough projectives, one has \(X \cong \colim_{i} Q_{i}\), realising \(X\) as the colimit of a finite diagram of projectives.
  We then have
  \begin{align*}
    \cat{C}(X, \hom{M,&\colim\displaylimits_{j}N_{j}})
                  \cong \cat{C}(\colim\displaylimits_{i}Q_{i},\hom{M,\colim\displaylimits_{j}N_{j}})\\
                & \cong \lim_{i}\cat{C}(Q_{i},\hom{M,\colim\displaylimits_{j}N_{j}})
                  \cong \lim_{i} \cat{M}(Q_{i} \lact M, \colim\displaylimits_{j} N_{j})\\
                & \cong \lim_{i}\colim\displaylimits_{j} \cat{M}(Q_{i} \lact M, N_{j})
                  \cong \lim_{i}\colim\displaylimits_{j} \cat{C}(Q_{i}, \hom{M,N_{j}})\\
                & \cong \lim_{i} \cat{C}(Q_{i}, \colim\displaylimits_{j} \hom{M,N_{j}})
                  \cong \cat{C}(\colim\displaylimits_{i} Q_{i}, \colim\displaylimits_{j} \hom{M,N_{j}})\\
                & \cong \cat{C}(X,\colim\displaylimits_{j} \hom{M,N_{j}}).\qedhere
  \end{align*}
\end{proof}

By oppositising \cref{Cprojectives}, we obtain the following result.

\begin{proposition}\label{Cinjectives}
  Assume that \(\cat{C}\) and \(\cat{M}\) have enough injectives.
  Let \(M\) be a coclosed object of \(\cat{M}\). Then \(\cohom{M,\blank}\) is left exact if and only if \(M\) is \(\cat{C}\)\hyp{}injective.
\end{proposition}

\begin{definition}
  Let \(\cat{C}\) be an abelian monoidal,
  and \(\cat{M}\) an abelian \(\cat{C}\)\hyp{}module category,
  with both having enough projectives.
  A \(\cat{C}\)\hyp{}projective object \(M\) is called a \emph{\(\cat{C}\)\hyp{}projective \(\cat{C}\)\hyp{}generator}
  if for any projective object \(Q\) in \(\cat{M}\)
  there exists a projective object \(P\) in \(\cat{C}\),
  such that \(Q\) is a direct summand of \(P \lact M\).

  Analogously, if \(\cat{C}\) and \(\cat{M}\) instead have enough injectives,
  a \(\cat{C}\)\hyp{}injective object \(M\) is called a \emph{\(\cat{C}\)\hyp{}injective \(\cat{C}\)\hyp{}cogenerator}
  if any injective object \(J\) of \(\cat{M}\) is a direct summand of an object of the form \(I \lact M\),
  for an injective object \(I\) of \(\cat{C}\).
\end{definition}

\begin{proposition}\label{ProjReflect}
  Assume that \(\cat{C}\) and \(\cat{M}\) have enough projectives.
  Let \(M \in \cat{M}\) be a closed \(\cat{C}\)\hyp{}projective object.
  If \(M\) is a \(\cat{C}\)\hyp{}generator, then \(\hom{M, \blank}\) reflects zero objects.

  If \(\proj{\cat{C}}\) is a Krull--Schmidt category
  and every indecomposable projective object of \(\cat{C}\) is the projective cover of a simple object,
  then we have that \(\hom{M, \blank}\) reflects zero objects
  if and only if
  \(M\) is a \(\cat{C}\)\hyp{}generator.
\end{proposition}
\begin{proof}
  Assume that \(M\) is a \(\cat{C}\)\hyp{}generator. Let \(N\) be a non-zero object of \(\cat{M}\) and let \(Q \longtwoheadrightarrow N\) be an epimorphism from a projective object in \(\cat{M}\). Let \(P\) be a projective object of \(\cat{C}\) such that \(Q\) is a direct summand of \(P \lact M\). Then \(\cat{M}(P \lact M, N)\) is non-zero, since \(\cat{M}(Q,N)\) is a direct summand thereof. Thus
  \[
    0 \neq \cat{M}(P \lact M, N) \cong \cat{C}(P,\hom{M,N}),
  \]
  showing that \(\hom{M,N}\) is not zero.

  For the latter statement, let \(Q'\) be an indecomposable projective object of \(\cat{M}\), and let \(S\) be its simple top. Since \(\hom{M,S} \neq 0\), there is some \(V \in \cat{M}\) such that \(\cat{C}(V,\hom{M,S}) \neq 0\). Let \(P' \longtwoheadrightarrow V\) be an epimorphism from a projective object in \(\cat{C}\). Then
  \[
    0\neq \cat{C}(V, \hom{M,S}) \simeq \cat{C}(P',\hom{M,S}) \simeq \cat{C}(P' \lact M, S).
  \]
  Since \(M\) is \(\cat{C}\)\hyp{}projective, the object \(P' \lact M\) is projective. Thus, \(\cat{C}(P' \lact M, S)\) being non-zero implies that \(Q'\) is a direct summand of \(P' \lact M\).
\end{proof}

More closely following the classical case,
one alternatively defines a closed \(M \in \cat{M}\) as a \(\cat{C}\)\hyp{}generator if \(\hom{M, \blank}\) is faithful,
see~\cite[Definition~2.21]{douglas19}.

Oppositisation of \cref{ProjReflect}
yields a similar variant in terms of coclosed and \(\cat{C}\)\hyp{}injective objects.

\begin{proposition}\label{InjReflect}
  Let \(\cat{C}\) and \(\cat{M}\) have enough injectives,
  and let \(M \in \cat{M}\) be a coclosed \(\cat{C}\)\hyp{}injective object.
  If \(M\) is a \(\cat{C}\)\hyp{}cogenerator, then \(\cohom{M, \blank}\) reflects zero objects.

  If \(\inj{\cat{C}}\) is a Krull--Schmidt category
  and every indecomposable injective object of \(\cat{C}\) is the injective hull of a simple object,
  then \(\cohom{M, \blank}\) reflects zero objects
  if and only if
  \(M\) is a \(\cat{C}\)\hyp{}cogenerator.
\end{proposition}

Finally, we give a brief account of internally projective and injective objects in a semisimple module category.

\begin{proposition}
  Let \(\cat{C}\) be a monoidal category with enough projectives, and let \(\cat{M}\) be a semisimple \(\cat{C}\)\hyp{}module category.%
  \marginnote{\textnormal{%
      Since \(\cat{M}\) is semisimple,
      any object of \(\cat{M}\) is \(\cat{C}\)\hyp{}projective and \(\cat{C}\)\hyp{}injective.}}
  An object \(M \in \cat{M}\) is a \(\cat{C}\)\hyp{}generator if and only if any simple object \(S \in \cat{M}\) is a direct summand of an object of the form \(P \lact M\), for some \(P \in \proj{\cat{C}}\).
\end{proposition}
\begin{proof}
  Recall that, in a semisimple abelian category, every object is projective.
  In particular, \(P \lact M\) is always projective, for \(M \in \cat{M}\) and \(P \in \proj{\cat{C}}\),
  hence every object in \(\cat{M}\) is \(\cat{C}\)\hyp{}projective.

  Now, let \(M \in \cat{M}\) be a \(\cat{C}\)\hyp{}generator.
  Since a simple object \(S \in \cat{M}\) is projective,
  it immediately follows
  from \(M\) being a \(\cat{C}\)\hyp{}generator
  that \(S\) is a direct summand of \(P \lact M\), for some \(P \in \proj{\cat{C}}\).

  Lastly, assume that every simple is a direct summand of \(P \lact M\),
  for some \(P \in \cat{C}\),
  and let \(Q \in \proj{\cat{M}}\).
  As \(\cat{M}\) is semisimple, we have \(Q \cong \oplus_i S_i\) for simples \(S_i\) in \(\cat{M}\).
  Since \(S_i\) is a direct summand of \(P_i \lact M\) for some \(P_{i} \in \cat{C}\), one calculates
  \[
    Q \,\cong\, \oplus_i S_i \,\subseteq_{\oplus}\, \oplus_i (P_i \lact M) \,\cong\, (\oplus_i P_i) \lact M \,\eqdef\, P \lact M.\qedhere
  \]
\end{proof}

\section{Reconstruction for lax module endofunctors}\label{sec:lax-module-reconstruction}

\textsc{The following is} our most general module categorical reconstruction result.

\begin{theorem}\label{mainnonsenseprojective}
  Let \(\cat{C}\) be a monoidal abelian category with enough projectives,
  \(\cat{M}\) an abelian \(\cat{C}\)\hyp{}module category with enough projectives,
  and assume that \(\ell \in \cat{M}\) is a closed \(\cat{C}\)\hyp{}projective \(\cat{C}\)\hyp{}generator.
  Then there is an equivalence of \(\cat{C}\)\hyp{}module categories
  \[
    \cat{M} \simeq \cat{C}^{\hom{\ell, \blank \,\lact\, \ell}},
  \]
  where \(\cat{C}^{\hom{\ell, \blank \,\lact\, \ell}}\) is endowed with the extended \(\cat{C}\)\hyp{}module structure
  by means of Linton coequalisers, see~\cref{extendingalways}.
\end{theorem}
\begin{proof}
  Following \cref{ex:regularend}, the functor \(\blank \lact \ell\) is a strong \(\cat{C}\)\hyp{}module functor.
  By \cref{porism:doctrinal-module-adjunctions}, its right adjoint \(\hom{\ell,\blank}\) is a lax \(\cat{C}\)\hyp{}module functor.
  Thus, the resulting monad \(\hom{\ell, \blank \lact \ell}\) is a right exact lax \(\cat{C}\)\hyp{}module monad.
  \cref{strongcomparisons} yields that the comparison functor \(\cat{C}_{\hom{\ell, \blank \,\lact\, \ell}} \to \cat{M}\) is a strong \(\cat{C}\)\hyp{}module functor.
  Furthermore, due to \cref{Cprojectives,ProjReflect}, \(\hom{\ell, \blank}\) is right exact and reflects zero objects,
  so we are able to apply \cref{thm:SmallBeck},
  whence the comparison functor \(\cat{M} \to \cat{C}^{\hom{\ell, \blank \,\lact\, \ell}}\) is an equivalence.
  Transporting the \(\cat{C}\)\hyp{}module structure along this equivalence,
  we obtain an extended \(\cat{C}\)\hyp{}module structure on \(\cat{C}^{\hom{\ell, \blank \,\lact\, \ell}}\),
  which,
  by \cref{thm:one-module-structure-on-EM},
  is necessarily that of \cref{extendingalways};
  this proves the result.
\end{proof}

Using \cref{extendingalways}, we can also formulate a converse.

\begin{theorem}\label{projcorrespondence}
  Let \(\cat{C}\) be a monoidal abelian category with enough projectives,
  and let \(T\) be a right exact lax \(\cat{C}\)\hyp{}module monad on \(\cat{C}\).
  Then \(T1\) is a closed \(\cat{C}\)\hyp{}projective \(\cat{C}\)\hyp{}generator in \(\cat{C}^T\),
  with the \(\cat{C}\)\hyp{}module category structure on the latter given by \cref{extendingalways}.
  There is a bijection
  \begin{align*}
    \faktor{\{\,(\cat{M},\ell) \text{ as in \cref{mainnonsenseprojective}}\,\}}{\cat{M} \simeq \cat{N}}
    &\xleftrightarrow{\ \cong\ }
      \Bigg\{
      \begin{aligned}
        &\text{Right exact lax \(\cat{C}\)\hyp{}module}\\
        &\text{monads on \(\cat{C}\)}
      \end{aligned}
      \Bigg\}
      /\, \cat{C}^T \!\simeq\! \cat{C}^S \\
    (\cat{M},\ell) &\longmapsto \hom{\ell,-\lact \ell} \\
    (\cat{C}^T, T1) &\longmapsfrom T
  \end{align*}
\end{theorem}
\begin{proof}
  The fact that \(T1\) is closed follows immediately from the isomorphisms
  \[
    \cat{C}^T(\blank \lact T1, \bblank)
    \overset{(i)}{\cong} \cat{C}^T(T(\blank \otimes 1), \bblank)
    \simeq \cat{C}(\blank \otimes 1, U^T(\bblank))
    \simeq \cat{C}(\blank, U^T(\bblank))
  \]
  showing that \(\hom{T1, \blank} \cong U^T\),
  where \((i)\) follows from \cref{actiononfree}.

  For \(P \in \proj{\cat{C}}\) projective, \(P \lact T(1) \cong TP\) is projective by \cref{finitemonads}.
  By the same result, \(T1\) is a \(\cat{C}\)\hyp{}projective generator,
  and \(\cat{C}^T\) has enough projectives.

  For the latter claim, notice that \(\cat{M} \simeq \cat{C}^{\hom{T1, \blank \,\lact\, T1}}\) by \cref{mainnonsenseprojective}.
  Lastly, \(\cat{C}^T \simeq \cat{C}^{\hom{T1, \blank \,\lact\, T1}}\) holds,
  since from \(\hom{T1, \blank} \cong U^T\) and, by uniqueness of adjoints, \(F^T \cong \blank \lact T1\)
  we deduce that \(T \cong \hom{T1, \blank \lact T1}\).
\end{proof}

Using \cref{Cinjectives,InjReflect} in place of \cref{Cprojectives,ProjReflect}, we find dual statements.

\begin{theorem}\label{mainnonsenseinjective}
  Let \(\cat{C}\) be a monoidal abelian category with enough injectives,
  \(\cat{M}\) an abelian \(\cat{C}\)\hyp{}module category with enough injectives,
  and assume that \(\ell \in \cat{M}\) be a coclosed \(\cat{C}\)\hyp{}injective \(\cat{C}\)\hyp{}cogenerator.

  Then there is an equivalence of \(\cat{C}\)\hyp{}module categories
  \[
    \cat{M} \simeq \cat{C}^{\cohom{\ell, \blank \,\lact\, \ell}},
  \]
  where \(\cat{C}^{\cohom{\ell, \blank \,\lact\, \ell}}\) is endowed with the extended \(\cat{C}\)\hyp{}module structure of \cref{extendingalways}.
\end{theorem}

\begin{theorem}\label{injcorrespondence}
  Let \(\cat{C}\) be a monoidal abelian category with enough injectives,
  and \(S\) a left exact oplax \(\cat{C}\)\hyp{}module comonad on \(\cat{C}\).
  Then \(\cat{C}^{S}\),
  endowed with the \(\cat{C}\)\hyp{}module category structure of \cref{extendingalways},
  is a \(\cat{C}\)\hyp{}module category,
  and \(S 1\) is a coclosed \(\cat{C}\)\hyp{}injective \(\cat{C}\)\hyp{}generator.
  There is a bijection
  \begin{align*}
    \faktor{\{\,(\cat{M}, \ell)\ \text{as in \cref{mainnonsenseinjective}}\,\}}{\cat{M} \simeq \cat{N}}
    &\leftrightarrow
    \left\{
      \begin{aligned}
        &\text{Left exact oplax}\ \cat{C}\text{-module}\\
        &\text{comonads on}\ \cat{C}
      \end{aligned}
      \right\}
      /\, \cat{C}^{S} \simeq \cat{C}^{G} \\
    (\cat{M},\ell) &\longmapsto \cohom{\ell, \blank \lact \ell} \\
    \big(\cat{C}^{S}, S 1\big) &\longmapsfrom S
  \end{align*}
\end{theorem}

Observe that \cref{mainnonsenseprojective,mainnonsenseinjective} do not make any finiteness assumptions on \(\cat{M}\)—%
such conditions are often imposed on it by the existence of a closed \(\cat{C}\)\hyp{}projective \(\cat{C}\)\hyp{}generator,
or a coclosed \(\cat{C}\)\hyp{}injective \(\cat{C}\)\hyp{}cogenerator.

\begin{proposition}
  Let \(\cat{C}\) be a finite abelian monoidal category
  and suppose \(\cat{M}\) is an abelian \(\cat{C}\)\hyp{}module category,
  such that there exists a closed \(\cat{C}\)\hyp{}projective \(\cat{C}\)\hyp{}generator \(\ell \in \cat{M}\).
  Then \(\cat{M}\) is finite abelian.
\end{proposition}
\begin{proof}
  By \cref{mainnonsenseprojective}, we have a \(\cat{C}\)\hyp{}module equivalence \(\cat{M} \simeq \cat{C}^{\hom{\ell,\blank \,\lact\, \ell}}\).
  Since \(\blank \lact \ell\) admits a right adjoint it is right exact.
  Since \(\ell\) is \(\cat{C}\)\hyp{}projective, the functor \(\hom{\ell, \blank}\) is also right exact.
  Thus \(\hom{\ell, \blank \lact \ell}\) is a right exact monad on the finite abelian category \(\cat{C}\).
  By \cref{finitemonads}, \(\cat{C}^{\hom{\ell, \blank \,\lact\, \ell}}\) is finite abelian.
\end{proof}

Using \cref{locomonads} in place of \cref{finitemonads}, one obtains the following.

\begin{proposition}
  Let \(\cat{C}\) be a locally finite abelian monoidal category,
  and let \(\cat{M}\) be an abelian \(\cat{C}\)\hyp{}module category with enough injectives,
  such that there exists a coclosed \(\Ind(\cat{C})\)\hyp{}injective \(\Ind(\cat{C})\)\hyp{}generator \(\ell \in \cat{M}\).
  Then the full subcategory of compact objects in \(\cat{M}\) is locally finite abelian.
\end{proposition}

In the presence of suitable finiteness conditions,
the characterisations of adjoint functors between finite and locally finite abelian categories
give sufficient conditions for an object in a module category to be closed or coclosed.

\begin{proposition}
  Let \(\cat{C}\) be a finite abelian monoidal category and let \(\cat{M}\) be a finite abelian \(\cat{C}\)\hyp{}module category.
  Then an object \(m \in \cat{M}\) is closed if and only if the functor \(\blank \lact m\) is right exact.
\end{proposition}
\begin{proof}
  This is an immediate consequence of \cref{finiteadjoints}.
\end{proof}

Similarly, using \cref{locallyfiniteadjoints}, one shows the following.

\begin{proposition}
  Let \(\cat{C}\) be a locally finite abelian monoidal category
  and \(\cat{M}\) a locally finite abelian \(\cat{C}\)\hyp{}module category.
  An object \(\ell \in \Ind(\cat{M})\) is coclosed with respect to the induced \(\Ind(\cat{C})\)\hyp{}module structure on \(\Ind(\cat{M})\)
  if and only if
  \(\blank \lact \ell \from \cat{C} \to \Ind(\cat{M})\) is left exact
  and the induced functor \(\blank \lact \ell\from \Ind(\cat{C}) \to \Ind(\cat{M})\) is quasi-finite.
\end{proposition}

\subsection{Rigid monoidal and (finite) tensor categories}

\textsc{We obtain an algebraic reconstruction result}
in case \(\cat{C}\) is rigid monoidal.

\begin{theorem}\label{thm:main-result-rigid-projective-case}
  Let \(\cat{C}\) be a rigid monoidal abelian category with enough projectives,
  \(\cat{M}\) an abelian \(\cat{C}\)\hyp{}module category with enough projectives,
  and assume that \(\ell \in \cat{M}\) is a closed \(\cat{C}\)\hyp{}projective \(\cat{C}\)\hyp{}generator.
  Then there is an algebra object \(A \in \cat{C}\) such that
  there is an equivalence of \(\cat{C}\)\hyp{}module categories
  \[
    \mathsf{mod}_{\cat{C}}A \simeq \cat{M}.
  \]
\end{theorem}
\begin{proof}
  By \cref{mainnonsenseprojective},
  we have \(\cat{C}^{\hom{\ell, \blank \,\lact\, \ell}} \simeq \cat{M}\)
  as \(\cat{C}\)\hyp{}module categories.
  By \cref{prop:lax-is-strong-when-rigid},
  the monad \(\hom{\ell, \blank \,\lact\, \ell}\from \cat{C}\to \cat{C}\) is a strong \(\cat{C}\)\hyp{}module functor,
  as it is lax and \(\cat{C}\) is rigid.
  The claim follows by \cref{prop:strong-module-monad-and-algebra-objects}.
\end{proof}

\begin{theorem}\label{thm:main-result-rigid-injective-case}
  Let \(\cat{C}\) be a rigid monoidal abelian category with enough injectives,
  \(\cat{M}\) an abelian \(\cat{C}\)\hyp{}module category,
  and assume that \(\ell \in \cat{M}\) is a coclosed \(\cat{C}\)\hyp{}injective \(\cat{C}\)\hyp{}cogenerator.
  Then there is a coalgebra object \(C \in \cat{C}\) such that there is an equivalence of
  \(\cat{C}\)\hyp{}module categories \(\mathsf{comod}_{\cat{C}}C \simeq \cat{M}\).
\end{theorem}

Recall the following result in case that \(\cat{C}\) is a multitensor category.

\begin{theorem}[{\cite[Theorem~7.10.1]{Etingof2015}}]
  Let \(\cat{C}\) be a finite multitensor category
  and let \(\cat{M}\) be an abelian \(\cat{C}\)\hyp{}module category, such that
  \begin{itemize}
    \item the functor \(\blank \lact \bblank\) is exact in the first variable
    (as \(\cat{C}\) is rigid, it is always exact in the second variable);
    \item there exists a \(\cat{C}\)\hyp{}projective \(\cat{C}\)\hyp{}generator.
  \end{itemize}
  Then there is an algebra object \(A\) in \(\cat{C}\) such that \(\cat{M} \simeq \mathsf{mod}_{\cat{C}}A\).
\end{theorem}

We generalise this statement to the locally finite setting.

\begin{theorem}\label{thm:multitensor-ind-reconstruction}
  Let \(\cat{C}\) be a multitensor category,
  and let \(\cat{M}\) be an abelian \(\cat{C}\)\hyp{}module category
  such that the \(\Ind(\cat{C})\)\hyp{}module category \(\Ind(\cat{M})\)
  admits a coclosed \(\Ind(\cat{C})\)\hyp{}injective \(\Ind(\cat{C})\)\hyp{}cogenerator \(\ell\).

  Then there is a coalgebra object \(C\) in \(\Ind(\cat{C})\)
  such that \(\Ind(\cat{M}) \simeq \mathsf{Comod}_{\Ind(\cat{C})}C\).
  Thus, \(\cat{M}\) is the category of compact \(C\)\hyp{}comodule objects.
\end{theorem}
\begin{proof}
  By \cref{mainnonsenseinjective}, we have \(\Ind(\cat{M}) \simeq {\Ind(\cat{C})}^{\cohom{\ell, \blank \,\lact\, \ell}}\)
  as \(\Ind(\cat{C})\)\hyp{}module categories.
  Observe that the \(\Ind(\cat{C})\)\hyp{}module category structure on both \(\Ind(\cat{C})\) and on \(\Ind(\cat{M})\)
  is the finitary extension of the respective \(\cat{C}\)\hyp{}module category structures,
  following \cref{cocompletingnonsense}.
  Similarly, \(\blank \lact \ell\from \Ind(\cat{C}) \to \Ind(\cat{M})\) is the extension of \(\blank \lact \ell\from \cat{C} \to \Ind(\cat{M})\).
  On the ind-completion, the left adjoint \(\cohom{\ell,\blank}\from \Ind(\cat{M}) \to \Ind(\cat{C})\),
  being finitary and preserving compact objects by \cref{leftadjointssmallobjects},
  restricts to an oplax \(\cat{C}\)\hyp{}module functor \(\cohom{\ell,\blank}\from \cat{M} \to \cat{C}\).

  Since \(\cat{C}\) is rigid, the restricted \(\cat{C}\)\hyp{}module functor \(\cohom{\ell,\blank}\from \cat{M} \to \cat{C}\)
  is in fact a strong \(\cat{C}\)\hyp{}module functor.
  Its finitary extension \(\cohom{\ell,\blank}\from \Ind(\cat{M}) \to \Ind(\cat{C})\) is a strong \(\Ind(\cat{C})\)\hyp{}module functor,
  and the comonad \(\cohom{\ell, \blank \lact \ell}\) is a strong \(\Ind(\cat{C})\)\hyp{}module monad.
  By \cref{prop:strong-module-monad-and-algebra-objects},
  it is of the form \(\blank \otimes C\) for a coalgebra object in \(\Ind(\cat{C})\),
  and hence
  \[
    \Ind(\cat{M}) \simeq {\Ind(\cat{C})}^{\cohom{\ell, \blank \,\lact\, \ell}} \simeq \mathsf{Comod}_{\Ind(\cat{C})}C.
  \]
  By \cref{locomonads}, \(\cat{M}\) consists of compact \(C\)\hyp{}comodule objects of \(\cat{C}\).
\end{proof}

\begin{corollary}\label{infinitecomodulecoalgebras}
  Let \(H\) be a Hopf algebra and let \(\cat{C} = \tetramodfd[H]{}\);
  in particular, we have a monoidal equivalence \(\Ind(\cat{C}) \simeq \Tetramod[H]{}\).
  Let \(\cat{M}\) be an abelian \(\cat{C}\)\hyp{}module category
  such that the \(\Ind(\cat{C})\)\hyp{}module category \(\Ind(\cat{M})\)
  admits a coclosed \(\Ind(\cat{C})\)\hyp{}injective \(\Ind(\cat{C})\)\hyp{}cogenerator.
  Then there is an \(H\)\hyp{}comodule coalgebra \(C\)
  such that there is an equivalence \(\Ind(\cat{M}) \simeq \mathsf{Comod}_{H}C\)
  of \(\Ind(\cat{C})\)\hyp{}module categories,
  restricting to a \(\cat{C}\)\hyp{}module equivalence \(\cat{M} \simeq \mathsf{comod}_{H}C\).
\end{corollary}

\section{An Eilenberg–Watts theorem for lax module monads}\label{Takyon}

\textsc{The formulation of the bijection}
of \cref{injcorrespondence} indicates a Morita aspect to the reconstruction theory for module categories,
as the equivalence relation imposed on monads strongly resembles—and can be specialised to—Morita equivalence of algebras.

In \cref{TurboEilenbergWatts},
we characterise the right exact lax \(\cat{C}\)\hyp{}module functors between Eilenberg--Moore categories for lax \(\cat{C}\)\hyp{}module monads
in terms of suitable bimodule objects in the category of endofunctors of \(\cat{C}\)%
—this gives a precise meaning to the notion of Morita equivalence of monads,
hence \cref{injectivebijection} yields a bijection
\begin{align*}
  \faktor{\{(\cat{M},\ell) \text{ as in \cref{mainnonsenseinjective}}\}}{\cat{M} \simeq \cat{N}}
  &\xleftrightarrow{\ \simeq\ }
    \left\{\begin{aligned}
      &\text{Left exact oplax \(\cat{C}\)\hyp{}mod-}\\
      &\text{ule comonads on \(\cat{C}\)}
    \end{aligned}
    \right\}\!\Big/\!\!\simeq_{\text{Morita}} \\
  (\cat{M},\ell) &\longmapsto \cohom{\ell,-\lact \ell} \\
  (\cat{M}^T, T1) &\longmapsfrom T
\end{align*}

Let us first introduce some vocabulary.

\begin{definition}
  Let \(\cat{A}\) be an abelian category and suppose that \(T\) and \(S\) are right exact monads on \(\cat{A}\).
  Viewing them as algebra objects in \(\mathsf{Rex}(\cat{A}, \cat{A})\),
  a \emph{\(T\)\hyp{}\(S\)\hyp{}biact functor} is an \(T\)\hyp{}\(S\)\hyp{}bimodule object in \(\mathsf{Rex}(\cat{A},\cat{A})\).
\end{definition}

In other words, a biact functor \(F\from \cat{A} \to \cat{A}\) is a triple \((F,\,F_{\mathsf{la}},\,F_{\mathsf{ra}})\),
consisting of a right exact endofunctor on \(\cat{A}\) together with transformations
\(F_{\mathsf{la}}\from T F \nt F\)
and
\(F_{\mathsf{ra}}\from F S \nt F\), satisfying the natural unitality and associativity axioms,
and commuting with each other in the sense that
\begin{equation}\label[diagram]{eq:comm-biact}
  \begin{mytikzcd}[ampersand replacement=\&]
    TFS \& TF \\
    FS \& F
    \arrow["{TF_{\mathsf{ra}}}", from=1-1, to=1-2]
    \arrow["{F_{\mathsf{la}}S}"', from=1-1, to=2-1]
    \arrow["{F_{\mathsf{la}}}", from=1-2, to=2-2]
    \arrow["{F_{\mathsf{ra}}}"', from=2-1, to=2-2]
  \end{mytikzcd}
\end{equation}
commutes.
The category of \(T\)\hyp{}\(S\)\hyp{}biact functors shall be denoted by \(\biact{T}{S}\).

\begin{hypothesis}
  For the rest of this section, let us fix two right exact monads \(T\) and \(S\) on an abelian category \(\cat{A}\).
\end{hypothesis}

\begin{remark}\label{factorthrough}
  For a \(T\)\hyp{}\(S\)\hyp{}biact functor \(F\),
  the left \(T\)\hyp{}action \(F_{\mathsf{la}}\) endows any object of the form \(Fa\), for \(a \in \cat{A}\),
  with the structure of a \(T\)\hyp{}module by defining \(\nabla_{Fa} \defeq F_{\mathsf{la}; a}\).
  Any morphism \(f \in \cat{A}(a,b)\) lifts to a map \(Ff\from Fa \to Fb\) of \(T\)\hyp{}modules
  by naturality of \(F_{\mathsf{la}}\).
  In particular, the functor \(F\) factors through the canonical forgetful functor \(U^T\from \cat{C}^T \to \cat{C}\).
  We write \(F = U^T \circ \tilde{F}\).
\end{remark}

\begin{definition}\label{def:biactewdef}
  Given a \(T\)\hyp{}\(S\)\hyp{}biact functor \(F\),
  define \(\tilde{F} \circ_{S} \blank\from \cat{A}^S \to \cat{A}^T\) by
  \begin{equation}\label{eq:eilenberg-watts-circ-coeq}
    \tilde{F} \circ_{S} \blank \defeq \coeq\Big(\!\!
    \begin{mytikzcd}[ampersand replacement=\&]
      {FS} \&\& {F}
      \arrow["{F \nabla}", shift left=2, from=1-1, to=1-3]
      \arrow["{F_{\mathrm{ra}}}"', shift right=2, from=1-1, to=1-3]
    \end{mytikzcd}\!\!\Big),
  \end{equation}
  where for \(x \in \cat{A}^S\) we endow \(Fx\) and \(FSx\)
  with the \(T\)\hyp{}module structures described in \cref{factorthrough}.
\end{definition}

The functoriality of \cref{def:biactewdef} follows from that of colimits.

\begin{remark}\label{rmk:leave-out-forgetful-functors}
  By naturality of \(F_{\mathsf{la}}\) and the commutativity condition of \cref{eq:comm-biact},
  both morphisms in \cref{eq:eilenberg-watts-circ-coeq} are morphisms in \(\cat{A}^T\);
  since the coequalisers in \(\cat{A}^{T}\) are created by \(U^T\),
  this coequaliser is one in \(\cat{A}^{T}\).
\end{remark}

\begin{proposition}\label{EilenbergWatts}
  Let \(T\) and \(S\) be right exact monads on an abelian category \(\cat{A}\).
  Writing \(\varepsilon^T \defeq \varepsilon^{F^T \dashv U^T}\) and \(\varepsilon^S \defeq \varepsilon^{F^S \dashv U^S}\),
  the following assignments extend to an equivalence of categories:%
  \marginnote{\textnormal{%
      Notice the similarity of this assignment and \cref{def: crossed_product_monad},
      only that we are now pushing \(\Phi\) along two adjunctions.
  }}
  \begin{align*}
    \mathsf{Rex}(\cat{A}^S,\cat{A}^{T}) &\longleftrightarrow \biact{T}{S} \\
    \Phi &\longmapsto \big(
        U^T \Phi F^S,\ %
        U^T \varepsilon^T \Phi F^S,\ %
        U^T \Phi \varepsilon^S F^S
        \big)\\
    \tilde{F} \circ_{S} \blank &\longmapsfrom F.
  \end{align*}
\end{proposition}
\begin{proof}
  Since \(F^S\) is right exact and \(U^T\) is exact, \(U^T \Phi F^S\) is right exact.

  It is easy to verify that the transformations \(U^S \varepsilon^T \Phi F^T\) and \(U^T \Phi \varepsilon^S F^S\) endow \(U^T \Phi F^S\) with the structure of a biact functor,
  and that for a natural transformation \(\alpha\from \Phi \nt \Phi'\),
  the resulting transformation \(U^T \alpha F^S\) is a morphism of biact functors.
  This proves the functoriality of \(\mathsf{Rex}(\cat{A}^{S},\cat{A}^{T}) \to \biact{T}{S}\).

  To see that the converse assignment is functorial,
  observe that given a morphism \(f\from F \nt G\) of biact functors,
  we obtain a morphism of forks in \(\cat{A}^{T}\):
  \[
    \begin{mytikzcd}[ampersand replacement=\&]
      {FSx} \&\& {Fx} \\
      \\
      {GSx} \&\& {Gx}
      \arrow["{F \nabla_{x}}", shift left=2, from=1-1, to=1-3]
      \arrow["{F_{\mathrm{ra}; x}}"', shift right=2, from=1-1, to=1-3]
      \arrow["{f_{Sx}}"', from=1-1, to=3-1]
      \arrow["{f_{x}}", from=1-3, to=3-3]
      \arrow["{G \nabla_{x}}", shift left=2, from=3-1, to=3-3]
      \arrow["{G_{\mathrm{ra};x}}"', shift right=2, from=3-1, to=3-3]
    \end{mytikzcd}
  \]
  This induces a morphism of coequalisers \(\tilde{F} \circ_{S} \blank \nt \tilde{G} \circ_{S} \blank\).

  We now prove that these functors define mutually quasi-inverse equivalences.
  First, observe that
  \(\widetilde{U^T \Phi F^S} = \Phi F^S \from \cat{A} \to \cat{A}^{T}\),
  since \(U^T\) is faithful and injective on objects,
  so it is a monomorphism in the 1-category \(\mathsf{Cat}_{\Bbbk}\).
  By slight abuse of notation, see \cref{rmk:leave-out-forgetful-functors},
  for \(x \in \cat{A}^{S}\) we have
  \[
    (\Phi F^S) \circ_{S} x = \coeq\Big(\!\!
    \begin{mytikzcd}[ampersand replacement=\&]
      {\Phi S^{2}x} \&\& {\Phi Sx}
      \arrow["{\Phi\mu^{S}_{x}}", shift left=2, from=1-1, to=1-3]
      \arrow["{\Phi S\nabla_{x}}"', shift right=2, from=1-1, to=1-3]
    \end{mytikzcd}\!\!\Big),
  \]
  which, since \(\Phi\) is right exact, is isomorphic to
  \[
    \Phi\Big(\coeq\Big(\!\!
    \begin{mytikzcd}[ampersand replacement=\&]
      {S^{2}x} \&\& {Sx}
      \arrow["{\mu^{S}_{x}}", shift left=2, from=1-1, to=1-3]
      \arrow["{S\nabla_{x}}"', shift right=2, from=1-1, to=1-3]
    \end{mytikzcd}\!\!\Big)\Big).
  \]
  Via the action \(\nabla_x\from Tx \to x\), the latter coequaliser is canonically isomorphic to \(x\),
  hence we obtain the following isomorphism, which is natural in \(\Phi\):
  \[
    \Phi \hat{\nabla}_{x}\from (\Phi F^S) \circ_{S} x \iso \Phi x.
  \]

  Conversely, let \(F\) be a biact functor and \(x \in \cat{A}\).
  We have
  \[
    \Big( U^T (\tilde{F} \circ_{S} \blank) F^S \Big)(x) =
    \coeq\Big(\!\!
    \begin{mytikzcd}[ampersand replacement=\&]
      {FS^{2}x} \&\& {FSx}
      \arrow["{F_{\mathsf{ra}; Sx}}", shift left=2, from=1-1, to=1-3]
      \arrow["{F\nabla_{Sx} = F\mu^S_{x}}"', shift right=2, from=1-1, to=1-3]
    \end{mytikzcd}
    \!\!\Big)
  \]
  Observe that \(F\), being a right \(S\)\hyp{}module object,
  is presented as a coequaliser in the category of such objects,
  which additionally splits in the underlying monoidal category \(\mathsf{Rex}(\cat{A},\cat{A})\).
  This is indicated in the following diagram:
  \[
    \begin{mytikzcd}[ampersand replacement=\&]
      {FS^{2}} \&\& {FS} \& {\coeq(F_{\mathsf{ra}} S, F \mu^S)} \\
      \&\&\& F
      \arrow["{F_{\mathsf{ra}} S}", shift left=2, from=1-1, to=1-3]
      \arrow["{F \mu^S}"', shift right=2, from=1-1, to=1-3]
      \arrow[two heads, from=1-3, to=1-4]
      \arrow["{F_{\mathsf{ra}}}"{description}, from=1-3, to=2-4]
      \arrow["{\exists! \Hat[mortarboard]{F}_{\mathsf{ra}}}", dashed, from=1-4, to=2-4]
      \arrow["\simeq"', dashed, from=1-4, to=2-4]
      \arrow["{F \eta^{S}}", curve={height=-12pt}, dotted, from=2-4, to=1-3]
    \end{mytikzcd}
  \]
  Since \(F_{\mathsf{ra}}\) is a morphism of left \(T\)\hyp{}module objects,
  we find that \(\Hat[mortarboard]{F}_{\mathsf{ra}}\) defines an isomorphism of biact functors
  \(U^T (\tilde{F} \circ_{S} \blank) F^S \cong F\)
  that is natural in \(F\).
\end{proof}

Let us now assume \(\cat{C}\) and \(\cat{M}\) to be abelian,
\(\blank \lact \bblank\) to be right exact in both variables,
and \(S\) and \(T\) to be right exact lax \(\cat{C}\)\hyp{}module monads on \(\cat{M}\).

\begin{definition}
  An \(S\)\hyp{}\(T\)\hyp{}bimodule object in the category \(\mathsf{RexLax}\cat{C}\mathsf{Mod}(\cat{M},\cat{M})\)
  of right exact lax \(\cat{C}\)\hyp{}module endofunctors of \(\cat{M}\) is called a
  \emph{lax \(\cat{C}\)\hyp{}module \(S\)\hyp{}\(T\)\hyp{}biact functor}.
  We denote the category of \(S\)\hyp{}\(T\)\hyp{}biact functors by \(\biactC{S}{T}\).
\end{definition}

In other words, a lax \(\cat{C}\)\hyp{}module biact functor is
a biact functor \((F,\,F_{\mathsf{la}},\,F_{\mathsf{ra}})\),
such that \(F\from \cat{M}\to \cat{M}\) is a right exact lax \(\cat{C}\)\hyp{}module functor
and \(F_{\mathsf{la}}\) and \(F_{\mathsf{ra}}\) are \(\cat{C}\)\hyp{}module transformations.

\begin{theorem}\label{ultraEilenbergWatts}
  The equivalence \(\mathsf{Rex}(\cat{M}^T,\cat{M}^{S}) \simeq \biact{S}{T}\)
  of \cref{EilenbergWatts} restricts to a faithful functor
  \[
    \cat{C}\mathsf{EW}\from \mathsf{LaxRex}(\cat{M}^{T}, \cat{M}^{S}) \to \biactC{S}{T},
  \]
  such that the following diagram commutes:
  \begin{equation}\label[diagram]{mikey}
    \begin{mytikzcd}[ampersand replacement=\&]
      {\mathsf{LaxRex}(\cat{M}^{T}, \cat{M}^{S})} \&\& {\biactC{S}{T}} \\
      {\mathsf{Rex}(\cat{M}^{T}, \cat{M}^{S})} \&\& {\biact{S}{T}}
      \arrow["\cat{C}\mathsf{EW}", from=1-1, to=1-3]
      \arrow[from=1-1, to=2-1]
      \arrow[from=1-3, to=2-3]
      \arrow["\simeq \text{ by~\ref{EilenbergWatts}}", from=2-1, to=2-3]
    \end{mytikzcd}
  \end{equation}
  The vertical arrows are the respective forgetful functors,
  and the categories \(\cat{M}^{T}\) and \(\cat{M}^{S}\) are endowed with the \(\cat{C}\)\hyp{}module structures of \cref{extendingalways}.
\end{theorem}
\begin{proof}
  By \cref{strongembeddingalways}, similarly to~\cite[Proposition~3.10]{aguiar18:monad},
  the functor \(F^T\from \cat{M} \to \cat{M}^{T}\) is a strong \(\cat{C}\)\hyp{}module functor, and so is \(F^S\).
  Thus, by \cref{porism:doctrinal-module-adjunctions}, the functor \(U^S\) is a lax \(\cat{C}\)\hyp{}module functor,
  and hence for a lax \(\cat{C}\)\hyp{}module functor \(\Phi\from \cat{M}^{T} \to \cat{M}^{S}\),
  the composite \(U^S \Phi F^T\) is a lax \(\cat{C}\)\hyp{}module functor.
  Further, the \(S\)\hyp{}\(T\)\hyp{}biact structure on \(U^S \Phi F^S\) given in \cref{EilenbergWatts}
  assembles to a lax \(\cat{C}\)\hyp{}module biact functor,
  and this extends also to \(\cat{C}\)\hyp{}module biact transformations arising from \(\cat{C}\)\hyp{}module transformations.

  This defines a functor \(\cat{C}\mathsf{EW}\), such that \cref{mikey} commutes.
  It is faithful, since so are the remaining functors in that diagram.
\end{proof}

\begin{lemma}\label{phigter}
  Let \(\Phi \in \mathsf{Rex}(\cat{M}^T,\cat{M}^S)\) and for all \(m \in \cat{M}\) suppose that
  \[
    {(\Phi T)}_{\mathsf{a};m}\from x \lact_{\cat{M}} \Phi Tm \to \Phi T(x \lact_{\cat{M}} m)
  \]
  is a lax \(\cat{C}\)\hyp{}module biact functor structure
  on the biact functor \(U^S \Phi F^T\).
  Then we have \({(\Phi T)}_{\mathsf{a}} = \cat{C}\mathsf{EW}(\Phi_{\mathsf{a};T(\blank)})\), where
  \[
    \Phi_{\mathsf{a};T(\blank)}\from x \lact_{\cat{M}^S} \Phi T(\blank) \nt \Phi(x \lact_{\cat{M}^T} T(\blank)) = \Phi(T(x\lact_{\cat{M}}\blank))
  \]
  is the unique natural transformation making the following diagram commute:
  \begin{equation}\label[diagram]{phiaeio}
    \begin{mytikzcd}[ampersand replacement=\&]
      {x \lact_{\cat{M}^S} \Phi T-} \\
      {S(x\lact_{\cat{M}}\Phi T-)} \& {S\Phi T x\lact_{\cat{M}} -} \& {\Phi T x\lact_{\cat{M}} \blank}
      \arrow["{\exists ! \Phi_{\mathsf{a}}}", dashed, from=1-1, to=2-3]
      \arrow[two heads, from=2-1, to=1-1]
      \arrow["{S{(\Phi T)}_{\mathsf{a}}}"', from=2-1, to=2-2]
      \arrow["{\Phi_{\mathsf{la}} Tx\lact\blank}"', from=2-2, to=2-3]
    \end{mytikzcd}
  \end{equation}
\end{lemma}
\begin{proof}
  In order to verify that \(\Phi_{\mathsf{a}}\) is well-defined,
  observe that the morphism \((\Phi_{\mathsf{la}}Tx\lact_{\cat{M}}\blank)\circ S{(\Phi T)}_{\mathsf{a}}\)
  coequalises \(SS_{\mathsf{a};x} \Phi T\) and \(Sx \lact_{\cat{M}} \Phi_{\mathsf{la}} T\):
  \[
    \tikzfig{phigter-1}
  \]
  Above, \((1)\) follows by \({(\Phi T)}_{\mathsf{la}}\) being a left \(S\)\hyp{}act structure,
  and \((2)\) by it being a \(\cat{C}\)\hyp{}module transformation.

  By the definition of the functor \(\cat{C}\mathsf{EW}\),
  we have
  \[
    \cat{C}\mathsf{EW}(\Phi_{\mathsf{a};T(\blank)}) = \Phi_{\mathsf{a};T(\blank)} \circ U^S_{\mathsf{a}; x, \Phi T(\blank)},
  \]
  where \(U^{S}_{\mathsf{a}}\) is the lax \(\cat{C}\)\hyp{}module functor structure on \(U^S\)
  induced by \cref{porism:doctrinal-module-adjunctions}.

  More generally, for \(m \in \cat{M}\) the map \(U^S_{\mathsf{a};x,m}\) is given by the composite
  \[
    x \lact_{\cat{M}} m \xrightarrow{\eta^S_{x \,\lact_{\!\cat{M}}\, m}} S(x \lact_{\cat{M}} m) \longtwoheadrightarrow x \lact_{\cat{M}^S} m,
  \]
  with the latter map being the projection onto the following coequaliser:
  \[
    \begin{mytikzcd}[ampersand replacement=\&]
      {\mathrm{coeq}\Big(T(x \lact T^2 m)} \& {T(x \lact  Tm)\Big)} \& {S(x \lact m)} \\
      {\mathrm{coeq}\Big(T(x \lact T^2 m)} \& {Tx \lact  Tm\Big)} \& {x \lact m}
      \arrow[shift left=2, from=1-1, to=1-2]
      \arrow[shift right=2, from=1-1, to=1-2]
      \arrow["{{T(x \lact \nabla_m)}}", from=1-1, to=2-1]
      \arrow["{{T(x  \lact \nabla_m)}}", from=1-2, to=2-2]
      \arrow[equals, from=1-3, to=1-2]
      \arrow[two heads, from=1-3, to=2-3]
      \arrow[shift left=2, from=2-1, to=2-2]
      \arrow[shift right=2, from=2-1, to=2-2]
      \arrow[equals, from=2-3, to=2-2]
    \end{mytikzcd}
  \]

  All in all, we find that \(\cat{C}\mathsf{EW}(\Phi_{\mathsf{a};T(\blank)})\)
  is given by the outer path of
  \[
    \begin{mytikzcd}[ampersand replacement=\&]
      \& {x \lact_{\cat{M}^S} \Phi T\blank} \\
      {S(x\lact_{\cat{M}}\Phi T\blank)} \& {S(x\lact_{\cat{M}}\Phi T\blank)} \& {S\Phi T x\lact_{\cat{M}} \blank} \& {\Phi T x\lact_{\cat{M}} \blank}
      \arrow["{\exists! \Phi_{\mathsf{a}}}", dashed, from=1-2, to=2-4]
      \arrow["{\eta_{x \,\lact_{\!\cat{M}}\, \Phi T\blank}}"', from=2-1, to=2-2]
      \arrow[two heads, from=2-2, to=1-2]
      \arrow["{S{(\Phi T)}_{\mathsf{a}}}"', from=2-2, to=2-3]
      \arrow["{\Phi_{\mathsf{la}} Tx\lact\blank}"', from=2-3, to=2-4]
    \end{mytikzcd}
  \]
  which, by the unitality of the \(S\)\hyp{}act structure \({(\Phi T)}_{\mathsf{la}}\), equals \({(\Phi T)}_{\mathsf{a}}\).
\end{proof}

For \(m \in \cat{M}^T\), let \(\Phi_{\mathsf{a};m}\) be the unique map such that
\begin{equation}\label[diagram]{tekken}
  \begin{mytikzcd}[ampersand replacement=\&]
    {x \lact_{\cat{M}^S} \Phi Tm} \&\& {x \lact_{\cat{M}^S} \Phi m} \\
    {\Phi T(x \lact m)} \\
    {\Phi(x \lact_{\cat{M}^T} Tm)} \&\& {\Phi(x \lact_{\cat{M}^T} m)}
    \arrow["{x \,\lact_{\cat{M}^S}\, \Phi \nabla_{m}}", two heads, from=1-1, to=1-3]
    \arrow["{\Phi_{\mathsf{a};Tm}}"', from=1-1, to=2-1]
    \arrow["{\exists! \Phi_{\mathsf{a};m}}", dashed, from=1-3, to=3-3]
    \arrow["{=}"', from=2-1, to=3-1]
    \arrow["{\Phi(x \,\lact_{\cat{M}^T}\, \nabla_{m})}"', two heads, from=3-1, to=3-3]
  \end{mytikzcd}
\end{equation}
commutes.
By the uniqueness in the defining property of \(\Phi_{\mathsf{a};m}\),
this defines a natural transformation
\[
  \Phi_{\mathsf{a}}\from \blank \lact_{\cat{M}^S} \Phi(\bblank) \nt \Phi(\blank \lact_{\cat{M}^T} \bblank).
\]

\begin{lemma}\label{bigdiagrams}
  For \(\Phi \in \mathsf{Rex}(\cat{M}^T,\cat{M}^S)\),
  the map \(\Phi_{\mathsf{a}}\from \blank \lact_{\cat{M}^S} \Phi(\bblank) \nt \Phi(\blank \lact_{\cat{M}^S} \bblank)\)
  extended from the morphisms \(\Phi_{\mathsf{a};Tm}\) of \cref{tekken}
  defines a lax \(\cat{C}\)\hyp{}module structure on \(\Phi\).
\end{lemma}
\begin{proof}
  In order to simplify notation,
  we write \(\lact\) for \(\lact_{\cat{M}}\)
  and \(\blact\) for \(\lact_{\cat{M}^T}\) and \(\lact_{\cat{M}^S}\),
  since each occurrence of either of these symbols is unambiguous with respect to which is used.
  We also write \(x \lact\) for \(x \lact \blank\),
  and suppress the horizontal composition symbols in \(\BiCat{Cat}_{\Bbbk}\),
  replacing them simply by concatenation.
  Hence, every occurrence of \(\circ\) in the diagrams that follow is a vertical composition of natural transformations.
  For example,
  \[
    S(y\lact)S(x\lact)\Phi T
    \xrightarrow{S(y\lact)(\Phi_{\mathsf{la}}T(x\lact) \circ S{(\Phi T)}_{\mathsf{a}, x})}
    S(y\lact)\Phi T(x \lact)
  \]
  represents the natural transformation
  \[
    S(y \lact S(x \lact \Phi T(\blank)))
    \xrightarrow{S(y \lact S{(\Phi T)}_{\mathsf{a}; y, x})}
    S(y \lact S \Phi T (x \lact \blank))
    \xrightarrow{S(y \lact \Phi_{\mathsf{la}}T(x \lact \blank))}
    S(y \lact \Phi T (x \lact \blank)).
  \]

  Since \(\Phi_{\mathsf{a}}\) is uniquely determined by \(\Phi_{\mathsf{a};T(\blank)}\),
  it suffices to show that
  \[
    \begin{mytikzcd}[ampersand replacement=\&,cramped]
      {y \blact \Phi(x \blact T)} \&\&\& {\Phi(y \blact x \blact T)} \\
      {y \blact x \blact \Phi T} \& {y \blact \Phi T (x\lact)} \& {\Phi T(y\lact)(x\lact)} \\
      {y \otimes x \blact \Phi T} \& {\Phi T (y \otimes x)\lact} \&\& {\Phi((y \otimes x) \blact T)}
      \arrow["{\Phi_{\mathsf{a}; y, x \blact T}}", from=1-1, to=1-4]
      \arrow["{\Phi\alpha_2}", from=1-4, to=3-4]
      \arrow["{y \blact \Phi_{\mathsf{a};x,T}}", from=2-1, to=1-1]
      \arrow["{{\alpha_{1}}}", from=2-1, to=2-2]
      \arrow["{{\alpha_{2}}}"', from=2-1, to=3-1]
      \arrow[Rightarrow, no head, from=2-2, to=1-1]
      \arrow["{{\alpha_{3}}}", from=2-2, to=2-3]
      \arrow[Rightarrow, no head, from=2-3, to=1-4]
      \arrow["{{\alpha_{5}}}"{description}, from=2-3, to=3-2]
      \arrow["{{\alpha_{4}}}", from=3-1, to=3-2]
      \arrow[Rightarrow, no head, from=3-2, to=3-4]
    \end{mytikzcd}
  \]
  commutes, where we define
  \(\alpha_{1} \defeq y \blact \Phi_{\mathsf{a},x}\),
  \(\alpha_{3} \defeq \Phi_{\mathsf{a},y}(x\lact)\),
  \(\alpha_{4} \defeq \Phi_{\mathsf{a},y\otimes x}\),
  and \(\alpha_{5} \defeq \Phi T(\lact_{\mathsf{a},y,x})\).
  The morphism \(\alpha_{2}\) is obtained from \cref{fig:def-of-alpha2},
  in which every marked epimorphism is the coequaliser of the pair preceding it,
  every dashed arrow is an induced morphism between coequalisers coming from a morphism of diagrams,
  and the dotted arrows come from the universal property of coequalisers.
  The notation for the coequaliser \((y,x) \blact \Phi T\) is
  to emphasise the connection with the multiactegorical approach of \cref{extendingeverything}.
  \begin{figure}[htbp]
    \centering
    \vspace{-\onelineskip}\mathscale{0.92}{\begin{mytikzcd}[ampersand replacement=\&,cramped]
        {S(y\lact)SS(x\lact)S\Phi T} \&\&\& {S(y\lact )SS(x\lact)\Phi T} \& {S(y\lact)S(x\blact \Phi T)} \\
        \\
        {S(y\lact )S(x\lact )S\Phi T} \&\&\& {S(y\lact )S(x\lact )\Phi T} \& {S(y\lact)(x \blact \Phi T)} \\
        \&\&\&\& {y\blact x \blact \Phi T} \\
        {S(y\lact )(x\lact )S\Phi T} \&\&\& {S(y\lact )(x\lact )\Phi T} \& {(y,x)\blact \Phi T} \\
        \\
        {S(y\otimes x\lact)S\Phi T} \&\&\& {(S(y\otimes x)\lact)\Phi T} \& {(y \otimes x) \blact \Phi T}
        \arrow["{{S(y\lact)S(\mu \circ SS_{\mathsf{a},x})\Phi T}}", shift left=2, from=1-1, to=1-4]
        \arrow["{{S(y\lact)SS(x\lact)\Phi_{\mathsf{la}} T}}"', shift right=2, from=1-1, to=1-4]
        \arrow["{{(\mu \circ S_{\mathsf{a},y})S(x\lact )S\Phi T}}"', shift right=3, from=1-1, to=3-1]
        \arrow["{{S(y\lact)\mu (x\lact )S\Phi T}}", shift left=3, from=1-1, to=3-1]
        \arrow[two heads, from=1-4, to=1-5]
        \arrow["{{(\mu\circ S_{\mathsf{a,y}})S(x\lact )\Phi T}}"', shift right=3, from=1-4, to=3-4]
        \arrow["{{S(y\lact)\mu (x\lact)\Phi T}}", shift left=3, from=1-4, to=3-4]
        \arrow[shift right=3, dashed, from=1-5, to=3-5]
        \arrow[shift left=3, dashed, from=1-5, to=3-5]
        \arrow["{{S(y\lact)(\mu \circ SS_{\mathsf{a,x}})\Phi T}}", shift left=2, from=3-1, to=3-4]
        \arrow["{{S(y\lact)S(x\lact)\Phi_{\mathsf{la}}T}}"', shift right=2, from=3-1, to=3-4]
        \arrow["{{(\mu \circ S_{\mathsf{a},y})(x\lact)S\Phi T}}"', two heads, from=3-1, to=5-1]
        \arrow[two heads, from=3-4, to=3-5]
        \arrow["{{(\mu\circ S_{\mathsf{a},y})(x\lact)\Phi T}}"', two heads, from=3-4, to=5-4]
        \arrow[two heads, from=3-5, to=4-5]
        \arrow["{{\exists! \beta_{0}}}"', dotted, from=4-5, to=5-5]
        \arrow["\cong", dotted, from=4-5, to=5-5]
        \arrow["{{\alpha_{2}}}", shift left=2, curve={height=-30pt}, dotted, from=4-5, to=7-5]
        \arrow["{{S(\mu \circ (S_{\mathsf{a},y}(x\lact)) \circ ((y\lact)S_{\mathsf{a},x}))\Phi T}}", shift left=2, from=5-1, to=5-4]
        \arrow["{{S(y\lact )(x\lact )\Phi_{\mathsf{la}} T}}"', shift right=2, from=5-1, to=5-4]
        \arrow["{{S(\lact_{\mathsf{a},y,x})S\Phi T}}"', from=5-1, to=7-1]
        \arrow[two heads, from=5-4, to=5-5]
        \arrow["{{S(\lact_{\mathsf{a},y,x})\Phi T}}", from=5-4, to=7-4]
        \arrow["{{\exists! \beta_{1}}}"', dotted, from=5-5, to=7-5]
        \arrow["\cong", dotted, from=5-5, to=7-5]
        \arrow["{{S(\mu\circ S_{\mathsf{a},y\otimes x})\Phi T}}", shift left=2, from=7-1, to=7-4]
        \arrow["{{S(y\otimes x\lact)\Phi_{\mathsf{la}} T}}"', shift right=2, from=7-1, to=7-4]
        \arrow[two heads, from=7-4, to=7-5]
      \end{mytikzcd}}\vspace{-\onelineskip}% XXX
    \scaption[-1cm]{%
      Definition of the morphism \(\alpha_{2}\).%
      \label{fig:def-of-alpha2}%
    }%
  \end{figure}

  We now have the diagram
  \[
    \scalebox{0.8}{\begin{mytikzcd}[ampersand replacement=\&, cramped, column sep=small]
	{S(y\lact )S(x\lact )\Phi T} \&\&\& {S(y\lact)\Phi T(x\lact)} \\
	\& {S(y\lact)(x\blact\Phi T)} \\
	{S(y\lact )(x\lact )\Phi T} \& {y\blact x \blact\Phi T} \&\& {y\blact\Phi T(x\lact)} \& {\Phi T (y\lact)(x\lact)} \\
	\& {(y \otimes x)\blact \Phi T} \\
	{S((y\otimes x)\lact )\Phi T} \&\&\&\& {\Phi T ((y\otimes x)\lact)}
	\arrow[""{name=0, anchor=center, inner sep=0}, "{S(y\lact)(\Phi_{\mathsf{la}}T(x\lact)\circ S{(\Phi T)}_{\mathsf{a},x})}", from=1-1, to=1-4]
	\arrow[two heads, from=1-1, to=2-2]
	\arrow[""{name=1, anchor=center, inner sep=0}, "{(\mu(y\lact) \circ S_{\mathsf{a},y}) (x\lact )\Phi T}"', from=1-1, to=3-1]
	\arrow[two heads, from=1-4, to=3-4]
	\arrow[""{name=2, anchor=center, inner sep=0}, "{((\Phi_{\mathsf{la}}(y\lact))\circ (S{(\Phi T)}_{\mathsf{a},y}))(x\lact)}", from=1-4, to=3-5]
	\arrow[""{name=3, anchor=center, inner sep=0}, "{\alpha'_{1}}"{description}, dashed, from=2-2, to=1-4]
	\arrow[two heads, from=2-2, to=3-2]
	\arrow["{S(\lact_{\mathsf{a},y,x})\Phi T}"', from=3-1, to=5-1]
	\arrow[""{name=4, anchor=center, inner sep=0}, "{\alpha_{1}}", from=3-2, to=3-4]
	\arrow[""{name=5, anchor=center, inner sep=0}, "{\alpha_{2}}", from=3-2, to=4-2]
	\arrow["{\alpha_{3}}", from=3-4, to=3-5]
	\arrow["{\alpha_{5}}"{description}, from=3-5, to=5-5]
	\arrow["{\alpha_{4}}"{description}, from=4-2, to=5-5]
	\arrow[two heads, from=5-1, to=4-2]
	\arrow[""{name=6, anchor=center, inner sep=0}, "{(\Phi_{\mathsf{la}}((y\otimes x)\lact))\circ (S{(\Phi T)}_{\mathsf{a},y\otimes x})}"', from=5-1, to=5-5]
	\arrow["{(2)}"{description}, draw=none, from=1, to=5]
	\arrow["{(1')}"{description}, draw=none, from=2-2, to=0]
	\arrow["{(1)}"{description}, draw=none, from=4, to=3]
	\arrow["{(3)}"{description}, draw=none, from=3-4, to=2]
	\arrow["{(4)}"{description}, draw=none, from=4-2, to=6]
      \end{mytikzcd}}
  \]
  whose labelled faces all commute,
  and where the morphism decorated by the label of the face it is part of
  is defined as that making said face commute,
  via the universal property of coequalisers.

  Our aim is to show the commutativity of the inner unlabelled face.
  Since all of its morphisms are defined by the remaining inner (commutative) faces,
  and we can reach \(y \blact x \blact \Phi T\) from \(S(y \lact)S(x \lact)\Phi T\) with two epimorphisms,
  this is implied by the commutativity of the outer face,
  which follows by the commutativity of the inner faces of \cref{fig:christmasmorning},
  where
  \begin{itemize}[noitemsep]
    \item faces (1), (3), and (4) commute by the interchange law in \(\BiCat{Cat}_{\Bbbk}\);
    \item face (2) commutes by \(\Phi_{\mathsf{la}}\from S \Phi T \nt \Phi T\) being a \(\cat{C}\)\hyp{}module map;
    \item face (5) commutes by the associativity of action \(\Phi_{\mathsf{la}}\);
    \item face (6) commutes by \(\Phi T\) being a lax \(\cat{C}\)\hyp{}module functor; and
    \item face (7) commutes by naturality of \(\Phi_{\mathsf{la}}\).\qedhere
  \end{itemize}
  \begin{figure}[htbp]
    \centering
    \scalebox{0.9}{\begin{mytikzcd}[ampersand replacement=\&,cramped,column sep=small, row sep=small]
	{S(y\lact )S(x\lact )\Phi T} \&\& {S(y\lact )S\Phi T(x\lact)} \&\&\&\& {S(y\lact)\Phi T(x\lact)} \\
	\\
	{SS(y\lact )(x\lact )\Phi T} \&\& {SS(y\lact )\Phi T(x\lact)} \&\& {SS\Phi T (y\lact)(x\lact)} \&\& {S\Phi T (y\lact)(x\lact)} \\
	\\
	{S(y\lact )(x\lact )\Phi T} \&\& {S(y\lact )\Phi T(x\lact)} \&\& {S\Phi T (y\lact)(x\lact)} \&\& {\Phi T (y\lact)(x\lact)} \\
	\\
	{S((y\otimes x)\lact )\Phi T} \&\&\&\& {S\Phi T ((y\otimes x)\lact)} \&\& {\Phi T ((y\otimes x)\lact)}
	\arrow[""{name=0, anchor=center, inner sep=0}, "\raisebox{0.6em}{\(S(y\lact )S{(\Phi T)}_{\mathsf{a},x}\)}", from=1-1, to=1-3]
	\arrow["{SS_{\mathsf{a},y} (x\lact )\Phi T}"', from=1-1, to=3-1]
	\arrow[""{name=1, anchor=center, inner sep=0}, "{S(y\lact)\Phi_{\mathsf{la}}T(x\lact)}", from=1-3, to=1-7]
	\arrow["{SS_{\mathsf{a},y}\Phi T(x\lact)}", from=1-3, to=3-3]
	\arrow["{S{(\Phi T)}_{\mathsf{a},y}(x\lact)}", from=1-7, to=3-7]
	\arrow[""{name=2, anchor=center, inner sep=0}, "\raisebox{0.6em}{\(SS(y\lact){(\Phi T)}_{\mathsf{a},x}\)}", from=3-1, to=3-3]
	\arrow["{\mu (y \lact)(x\lact)\Phi T}"', from=3-1, to=5-1]
	\arrow[""{name=3, anchor=center, inner sep=0}, "\raisebox{-0.8em}{\(SS{(\Phi T)}_{\mathsf{a},y}(x \lact )\)}"', from=3-3, to=3-5]
	\arrow["{\mu (y\lact)\Phi T(x\lact)}"', from=3-3, to=5-3]
	\arrow[""{name=4, anchor=center, inner sep=0}, "\raisebox{-0.8em}{\(S\Phi_{\mathsf{la}}T(y\lact)(x\lact)\)}"', from=3-5, to=3-7]
	\arrow["{\mu\Phi T (y\lact)(x\lact)}"', from=3-5, to=5-5]
	\arrow["{\Phi_{\mathsf{la}}T(y\lact)(x\lact)}", from=3-7, to=5-7]
	\arrow[""{name=5, anchor=center, inner sep=0}, "{S(y\lact){(\Phi T)}_{\mathsf{a},x}}"', from=5-1, to=5-3]
	\arrow["{S(\lact_{\mathsf{a},y,x})\Phi T}"', from=5-1, to=7-1]
	\arrow[""{name=6, anchor=center, inner sep=0}, "{S{(\Phi T)}_{\mathsf{a},y}(x\lact)}"', from=5-3, to=5-5]
	\arrow[""{name=7, anchor=center, inner sep=0}, "{\Phi_{\mathsf{la}} T (y\lact)(x\lact)}"', from=5-5, to=5-7]
	\arrow["{S\Phi T(\lact_{\mathsf{a},y,x})}"', from=5-5, to=7-5]
	\arrow["{\Phi T(\lact_{\mathsf{a},y,x})}", from=5-7, to=7-7]
	\arrow[""{name=8, anchor=center, inner sep=0}, "{S{(\Phi T)}_{\mathsf{a},y\otimes x}}"', from=7-1, to=7-5]
	\arrow[""{name=9, anchor=center, inner sep=0}, "\raisebox{-0.8em}{\(\Phi_{\mathsf{la}}T((y\otimes x)\lact)\)}"', from=7-5, to=7-7]
	\arrow["{(1)}"{description,pos=0.7}, shift left=2, draw=none, from=2, to=0]
	\arrow["{(2)}"{description}, draw=none, from=1, to=3-5]
	\arrow["{(3)}"{description}, shift right=5, draw=none, from=2, to=5]
	\arrow["{(4)}"{description,pos=0.6}, shift right=5, draw=none, from=3, to=6]
	\arrow["{(5)}"{description}, draw=none, from=4, to=7]
	\arrow["{(6)}"{description}, draw=none, from=5-3, to=8]
	\arrow["{(7)}"{description}, draw=none, from=7, to=9]
      \end{mytikzcd}}\vspace{-0.8cm}% XXX
    \scaption[-4cm]{%
      The arrow \(\Phi_{\mathsf{a}}\) defines a lax \(\cat{C}\)\hyp{}module structure on \(\Phi\).%
      \label{fig:christmasmorning}%
    }
  \end{figure}
\end{proof}

\begin{lemma}\label{jupiter}
  Let \(\Phi, \Phi' \in \mathsf{LaxRex}(\cat{M}^T,\cat{M}^S)\)
  and \(\phi\from U^S \Phi F^T \nt U^S \Phi' F^T\) be a lax \(\cat{C}\)\hyp{}module biact transformation.
  Then
  \[
    \varphi \defeq \phi \circ_{T} \blank \from \Phi \nt \Phi'
  \]
  is a \(\cat{C}\)\hyp{}module transformation.
\end{lemma}
\begin{proof}
  First, observe that \(\varphi\) is determined by \(\varphi_{T(\blank)}\), as indicated by the diagram
  \[
    \begin{mytikzcd}[ampersand replacement=\&]
      {\Phi T^{2}m} \& {\Phi Tm} \& {\Phi m} \\
      {\Phi' T^{2}m} \& {\Phi' Tm} \& {\Phi' m}
      \arrow["{\Phi\mu_{m}}"', shift right, from=1-1, to=1-2]
      \arrow["{\Phi T\nabla_{m}}", shift left, from=1-1, to=1-2]
      \arrow["{\varphi_{T^2m}}"', from=1-1, to=2-1]
      \arrow[two heads, from=1-2, to=1-3]
      \arrow["{\varphi_{Tm}}", from=1-2, to=2-2]
      \arrow["{\exists! \varphi_{m}}", dashed, from=1-3, to=2-3]
      \arrow["{\Phi' T\nabla_{m}}", shift left, from=2-1, to=2-2]
      \arrow["{\Phi'\mu_{m}}"', shift right, from=2-1, to=2-2]
      \arrow[two heads, from=2-2, to=2-3]
    \end{mytikzcd}
  \]
  Thus, it suffices to show that for any \(m \in \cat{M}\) and \(x \in \cat{C}\), we have that
  \begin{equation}\label{laxplax}
    \Phi'_{\mathsf{a};x,Tm} \circ (x \blact \phi_{m}) = \phi_{x\lact m} \circ \Phi_{\mathsf{a},x,T(m)}.
  \end{equation}

  Consider the following diagram:
  \[
    \begin{mytikzcd}[ampersand replacement=\&]
      {S(x\lact \Phi T)} \&\&\& {S\Phi Tx\lact \blank} \\
      \& {x \blact \Phi T} \&\&\& {\Phi Tx\lact \blank} \\
      {S(x\lact \Phi' T)} \&\&\& {S\Phi' Tx\lact \blank} \\
      \& {x \blact \Phi' T} \&\&\& {\Phi' Tx\lact \blank}
      \arrow["{S{(\Phi T)}_{\mathsf{a},x}}", from=1-1, to=1-4]
      \arrow[two heads, from=1-1, to=2-2]
      \arrow["{S(x\lact \phi)}"', from=1-1, to=3-1]
      \arrow["{\Phi_{\mathsf{la}}Tx\lact \blank}", from=1-4, to=2-5]
      \arrow["{S\phi'x\lact\blank}"{pos=0.8}, from=1-4, to=3-4]
      \arrow[crossing over, "{\Phi_{\mathsf{a},x,T(\blank)}}"{pos=0.4}, from=2-2, to=2-5]
      \arrow["{\phi_{x\lact \blank}}", from=2-5, to=4-5]
      \arrow["{S{(\Phi' T)}_{\mathsf{a},x}}"{pos=0.7}, from=3-1, to=3-4]
      \arrow[crossing over, "{x \blact\phi_{\blank}}"'{pos=0.3}, from=2-2, to=4-2]
      \arrow[two heads, from=3-1, to=4-2]
      \arrow["{\Phi'_{\mathsf{la}}Tx\lact \blank}", from=3-4, to=4-5]
      \arrow["{\Phi'_{\mathsf{a},x,T(\blank)}}"', from=4-2, to=4-5]
    \end{mytikzcd}
  \]
  Its top and bottom faces commute by definition of \(\Phi_{\mathsf{a}}\),
  as seen from \cref{phiaeio};
  its left face commutes by the definition of \(x \blact \blank\);
  and its right face by \(\phi\) being a biact transformation.
  The back face commutes since \(\phi\) is a \(\cat{C}\)\hyp{}module transformation.
  The commutativity of the front face is precisely \cref{laxplax},
  and since the top-left projection map is an epimorphism,
  it suffices to verify its commutativity after precomposing with it,
  which follows easily from the commutativity of the remaining faces.
\end{proof}

\begin{theorem}\label{TurboEilenbergWatts}
  The faithful functor
  \[
    \cat{C}\mathsf{EW}\from \mathsf{LaxRex}(\cat{M}^T, \cat{M}^S) \to \biactC[\cat{C}]{S}{T}
  \]
  of \cref{ultraEilenbergWatts} is an equivalence of categories.
\end{theorem}
\begin{proof}
  Fullness follows from \cref{jupiter}, so it is left to prove that \(\cat{C}\mathsf{EW}\) is essentially surjective.
  Suppose that \(F \in \biactC[\cat{C}]{T}{S}\).
  Since the functor of \cref{EilenbergWatts} is essentially surjective,
  we may assume that \(F \cong U^S \Phi F^T\),
  for some right exact \(\Phi\from \cat{M}^T \to \cat{M}^S\).
  By \cref{bigdiagrams}, \((\Phi, \Phi_{\mathsf{a}}) \in  \mathsf{LaxRex}(\cat{M}^T,\cat{M}^S)\),
  and \(F \cong \cat{C}\mathsf{EW}((\Phi, \Phi_{\mathsf{a}}))\) due to \cref{phigter},
  giving essential surjectivity.
\end{proof}

\begin{definition}
  We say that two right exact lax \(\cat{C}\)\hyp{}module monads \(T\) and \(S\) on \(\cat{M}\)
  are \emph{Morita equivalent} if there are
  \(F \in \biactC[\cat{C}]{T}{S}\) and
  \(G \in \biactC[\cat{C}]{S}{T}\)
  such that \(G \circ_{T} F \cong \Id_{\cat{M}^T}\)
  and \(F \circ_{S} G \cong \Id_{\cat{M}^{S}}\)
  as lax \(\cat{C}\)\hyp{}module biact functors.
\end{definition}

\begin{proposition}\label{Moritabybimodules}
  Two right exact lax \(\cat{C}\)\hyp{}module monads \(T\) and \(S\) on \(\cat{M}\)
  are Morita equivalent if and only if
  there is a \(\cat{C}\)\hyp{}module equivalence \(\cat{M}^T \simeq \cat{M}^{S}\),
  where the Eilenberg--Moore categories are endowed with the extended Linton coequaliser structure of \cref{extendingalways}.
\end{proposition}
\begin{proof}
  This is a direct consequence of \cref{TurboEilenbergWatts},
  using that an equivalence of categories is right exact.
\end{proof}

Combining \cref{Moritabybimodules} with \cref{injcorrespondence}, we find the following.

\begin{theorem}\label{injectivebijection}
  Let \(\cat{C}\) be a monoidal abelian category with enough injectives,
  and let \(T\) be a left exact lax \(\cat{C}\)\hyp{}module comonad on \(\cat{C}\).
  There is a bijection
  \begin{align*}
    \faktor{\{\,(\cat{M},\ell) \text{ as in \cref{mainnonsenseinjective}}\,\}}{\cat{M} \simeq \cat{N}}
    &\xleftrightarrow{\ \simeq\ }
      \left\{\begin{aligned}
        &\text{Left exact oplax \(\cat{C}\)\hyp{}module}\\
        &\text{comonads on \(\cat{C}\)}
      \end{aligned}
      \right\}\!\Big/\!\!\simeq_{\text{Morita}} \\
    (\cat{M},\ell) &\longmapsto \cohom{\ell,-\lact \ell} \\
    (\cat{M}^T, T1) &\longmapsfrom T
  \end{align*}
\end{theorem}

%%% Local Variables:
%%% mode: LaTeX
%%% TeX-master: "../main"
%%% TeX-command-extra-options: "-shell-escape -fmt=prec.fmt -file-line-error -synctex=1"
%%% End:
